\documentclass[12pt,a4paper]{article}
\usepackage{fontspec}
\usepackage{geometry}
\usepackage{enumitem}
\usepackage{fancyhdr}
\usepackage{lastpage}
\usepackage{hyperref}

\setmainfont{Noto Sans CJK JP}[Scale=1.0]

\geometry{margin=2.5cm}
\pagestyle{fancy}
\fancyhf{}
\rhead{Page \thepage\ of \pageref{LastPage}}
\lhead{BTC101 - Quiz (ja)}
\renewcommand{\headrulewidth}{0.4pt}

\title{Quiz: BTC101}
\author{Language: ja}
\date{\today}

\begin{document}

\maketitle
\thispagestyle{fancy}

\section*{Instructions}
Answer all questions by selecting the best option (A, B, C, or D). The answer key is provided at the end of this document.

\bigskip


\subsection*{Question 1}
ブロックサイズ戦争はいつ起こりましたか？

\begin{enumerate}[label=\Alph*., leftmargin=*]
\item 2008年に。
\item 2019年に。
\item 2017年に。
\item 2010年に。
\end{enumerate}

\bigskip


\subsection*{Question 2}
このビットコインコースの目的は何ですか？

\begin{enumerate}[label=\Alph*., leftmargin=*]
\item ビットコインのマイニング方法を紹介すること。
\item 新しいビットコインプロトコルの作成方法を紹介すること。
\item ビットコインの一般的な金融面の側面、その使用方法、および将来性について紹介すること。
\item ビットコインの規制方法を紹介すること。
\end{enumerate}

\bigskip


\subsection*{Question 3}
Bitcoinの二つの主な特徴は何ですか？

\begin{enumerate}[label=\Alph*., leftmargin=*]
\item Bitcoinは、株式の一種であり、債券の一種でもあります。
\item Bitcoinは、物理通貨であり、デジタルウォレットでもあります。
\item Bitcoinは、コンピュータプロトコルであり、かつ通貨単位でもあります。
\item Bitcoinは、中央集権的であり、かつ分散型通貨でもあります。
\end{enumerate}

\bigskip


\subsection*{Question 4}
このビットコイン講座の主な焦点は何ですか？

\begin{enumerate}[label=\Alph*., leftmargin=*]
\item ビットコインの金融面。
\item ビットコインの環境への影響。
\item ビットコインの規制面。
\item ビットコインのコーディング面。
\end{enumerate}

\bigskip


\subsection*{Question 5}
ビットコインが中立的な通貨であるとはどういう意味ですか？

\begin{enumerate}[label=\Alph*., leftmargin=*]
\item 政府によって規制されている。
\item どのような団体や機関にも制御されていない。
\item 一つの団体や機関によって制御されている。
\item 大きな取引にのみ使用される。
\end{enumerate}

\bigskip


\subsection*{Question 6}
サイファーパンクスとは誰ですか？

\begin{enumerate}[label=\Alph*., leftmargin=*]
\item デジタル時代におけるプライバシーと自由の役割を深く問い直し始めた人々のグループです。
\item 技術の使用に反対した人々のグループ。
\item 政府から情報を盗んだハッカーのグループ。
\item ビットコインやその他の暗号通貨を発明した人々のグループ。
\end{enumerate}

\bigskip


\subsection*{Question 7}
Cypherpunkマニフェストとは何ですか？

\begin{enumerate}[label=\Alph*., leftmargin=*]
\item Julian Assangeによって書かれた、技術の未来についての本。
\item プライバシーは基本的な権利であると主張する、Eric Hughesによって書かれたテキストです。
\item Bitcoinの原則を概説した文書。
\item デジタル通貨の重要性についてのDavid Chaumによって書かれたマニフェスト。
\end{enumerate}

\bigskip


\subsection*{Question 8}
DigiCashプロジェクトとは何でしたか？

\begin{enumerate}[label=\Alph*., leftmargin=*]
\item David Chaumによる匿名電子マネーを作成する試み。
\item 連邦準備制度による全ての現金をデジタル化する試み。
\item Bitcoinのためのデジタルウォレットを作成するプロジェクト。
\item Bitcoinに基づいた暗号通貨を作成するプロジェクト。
\end{enumerate}

\bigskip


\subsection*{Question 9}
サイバースペースの独立宣言とは何ですか？

\begin{enumerate}[label=\Alph*., leftmargin=*]
\item サイバースペースは政府の規制によって制御されるべきだと主張するテキスト。
\item いかなる規制からも海賊運動の独立を証明する宣言。
\item サイバースペースは物理的な領域とは異なる独自の領域であり、同じ法律の適用を受けるべきではないと主張するテキストです。
\item 伝統的な金融システムからのビットコインの独立を宣言する文書。
\end{enumerate}

\bigskip


\subsection*{Question 10}
Wei Daiのb-moneyとは何ですか？

\begin{enumerate}[label=\Alph*., leftmargin=*]
\item 匿名のデジタル通貨のアイデアであり、詐欺の検出は中央機関ではなく、評価者のコミュニティによって行われました。
\item CIAによって作成され、軍によって使用されたデジタルマネーの形態です。
\item Bitcoinより前に広く使用されていたデジタル通貨で、サイファーパンクスのグループによって管理されていました。
\item Bitcoinに対抗してWei Daiによって作成され、よりスケーラブルであることを目指したデジタル通貨です。
\end{enumerate}

\bigskip


\subsection*{Question 11}
ビットコインの技術を超えて何を象徴していますか？

\begin{enumerate}[label=\Alph*., leftmargin=*]
\item 伝統的な通貨よりも効率的な新しい形のデジタル通貨です。なぜなら、それはより速いからです。
\item 伝統的な金融システムに対する革命、そして透明性、分散化、個人の主権に基づく代替案です。
\item 個人がマイニングを通じてお金を稼ぐ方法です。
\item 犯罪者が違法取引を行うためのツールです。
\end{enumerate}

\bigskip


\subsection*{Question 12}
最初の通貨形態は何でしたか？

\begin{enumerate}[label=\Alph*., leftmargin=*]
\item 穀物や家畜などの商品。
\item 紙幣。
\item 金や銀。
\item デジタル資産。
\end{enumerate}

\bigskip


\subsection*{Question 13}
お金の三つの基本的な機能とは何ですか？

\begin{enumerate}[label=\Alph*., leftmargin=*]
\item 価値の保存、交換の媒体、計算の単位。
\item 価値の保存、生産の手段、計算の単位。
\item 交換の媒体、生産の手段、計算の単位。
\item 価値の保存、交換の媒体、生産の手段。
\end{enumerate}

\bigskip


\subsection*{Question 14}
効果的な通貨の主な特徴は何ですか？

\begin{enumerate}[label=\Alph*., leftmargin=*]
\item 非代替可能、分割可能、流動性がある。
\item 代替可能、分割可能、流動性がある。
\item 代替可能、分割可能、腐敗しやすい。
\item 代替可能、不可分、流動性がある。
\end{enumerate}

\bigskip


\subsection*{Question 15}
金をお金の形態として使用する際の欠点の一つは何ですか？

\begin{enumerate}[label=\Alph*., leftmargin=*]
\item 時間が経つと侵食される。
\item 容易に分割できない。
\item 生産が容易である。
\item 空間を越えてその価値を保持しない。
\end{enumerate}

\bigskip


\subsection*{Question 16}
なぜ金が標準的な交換媒体となったのか？

\begin{enumerate}[label=\Alph*., leftmargin=*]
\item その希少性、腐敗性、および分割可能性のため。
\item その豊富さ、腐敗性、および不可分性のため。
\item その豊富さ、耐久性、および不可分性のため。
\item その希少性、耐久性、および分割可能性のため。
\end{enumerate}

\bigskip


\subsection*{Question 17}
お金がコミュニケーションツールとして果たす役割は何ですか？

\begin{enumerate}[label=\Alph*., leftmargin=*]
\item 空間と時間を超えた価値の伝達を可能にし、相手方を信頼する必要がない。
\item 空間と時間を超えた価値の伝達を可能にし、相手方を信頼する必要がある。
\item 空間を超えた価値の伝達のみを可能にし、相手方を信頼する必要がない。
\item 時間を超えた価値の伝達のみを可能にし、相手方を信頼する必要がない。
\end{enumerate}

\bigskip


\subsection*{Question 18}
フィアット通貨の主な欠点は何ですか？

\begin{enumerate}[label=\Alph*., leftmargin=*]
\item その価値は通貨インフレーションによって常に高まります。
\item それらは簡単に運べるものではありません。
\item それらは分割できません。
\item その価値は通貨インフレーションによって常に侵食されます。
\end{enumerate}

\bigskip


\subsection*{Question 19}
Bitcoinが他の通貨形態に対して持つ主な利点は何ですか？

\begin{enumerate}[label=\Alph*., leftmargin=*]
\item それは分割不可能であり、運搬もできません。
\item 世界中のすべての商人によって商業で広く受け入れられています。
\item それは通貨インフレを受けやすい優れた例です。
\item 厳格に限定された供給量と中立性により、優れた価値の保存手段を提供します。
\end{enumerate}

\bigskip


\subsection*{Question 20}
なぜ金はグローバル化されたデジタルの世界でうまく機能しないのか？

\begin{enumerate}[label=\Alph*., leftmargin=*]
\item 分割や容易な交換を可能にする中央の実体が必要だからです。
\item 時間とともに腐食するからです。
\item 十分に希少ではないからです。
\item それは十分に価値がないからです。
\end{enumerate}

\bigskip


\subsection*{Question 21}
なぜビットコインは今日、広く商業で受け入れられていないのですか？

\begin{enumerate}[label=\Alph*., leftmargin=*]
\item 個人による採用が広くないためです。
\item それは小さい単位に分割できないためです。
\item それは輸送可能ではないためです。
\item それは流動性がないためです。
\end{enumerate}

\bigskip


\subsection*{Question 22}
今日の通貨の進化の道筋は何ですか？

\begin{enumerate}[label=\Alph*., leftmargin=*]
\item 生の石 -> コイン -> 紙幣 -> 銀行カード -> ブロックチェーン -> ライトニングネットワーク
\item 生の石 -> コイン -> 紙幣 -> 銀行カード -> ライトニングネットワーク
\item 生の石 -> コイン -> 銀行カード -> 紙幣 -> ブロックチェーン -> ライトニングネットワーク
\item 生の石 -> コイン -> 紙幣 -> ブロックチェーン -> ライトニングネットワーク
\end{enumerate}

\bigskip


\subsection*{Question 23}
なぜ金は持ち運びに適した素材ではないのですか？

\begin{enumerate}[label=\Alph*., leftmargin=*]
\item 室温で液体になるからです。
\item 小さな形に簡単に形成することができないからです。
\item それは密度が高く重いため、大量に持ち運ぶのが煩わしいからです。
\item それは移動させるには価値が高すぎるからです。
\end{enumerate}

\bigskip


\subsection*{Question 24}
フィアット通貨の主な利点は何ですか？

\begin{enumerate}[label=\Alph*., leftmargin=*]
\item 時間を超えて価値を保持するのに効果的です。
\item 中央の当事者によって制御されています。
\item 簡単に運べるわけではありません。
\item 簡単に分割できることです。
\end{enumerate}

\bigskip


\subsection*{Question 25}
次のうち、ビットコインが通貨として持たない利点はどれですか？

\begin{enumerate}[label=\Alph*., leftmargin=*]
\item 税金の支払いに使用されます。
\item 分割可能です。
\item 許可が不要です。
\item 持ち運びが可能です。
\end{enumerate}

\bigskip


\subsection*{Question 26}
健全な通貨はどのようにして作られるのか？

\begin{enumerate}[label=\Alph*., leftmargin=*]
\item 健全な通貨は、政府が通貨法を発行することによって強制される必要があります。
\item 健全な通貨は市場と自発的な合意から自然に生まれなければなりません。
\item 健全な通貨は、企業が暴力を用いることによって強制されなければなりません。
\item 健全な通貨は作られることはできません。
\end{enumerate}

\bigskip


\subsection*{Question 27}
ビットコインが新しい形のお金とみなされる理由は何ですか？

\begin{enumerate}[label=\Alph*., leftmargin=*]
\item 金に裏打ちされたデジタル通貨であるため。
\item 銀行のブロックチェーンが必要な通貨であるため。
\item 中央の組織がないデジタルマネーであるためです。
\item 連邦準備制度がそう言っているから。
\end{enumerate}

\bigskip


\subsection*{Question 28}
金が通貨としての主な利点は何ですか？

\begin{enumerate}[label=\Alph*., leftmargin=*]
\item 簡単に分割したり、長距離を運ぶことができます。
\item 確認が容易です。
\item 政府による規制がありません。
\item 非常に耐久性があることです。
\end{enumerate}

\bigskip


\subsection*{Question 29}
「フィアットマネー」の意味は何ですか？

\begin{enumerate}[label=\Alph*., leftmargin=*]
\item フィアットマネーは、政府がその価値を維持し、人々がその価値を信じているために価値がある通貨であり、金のような物理的な商品によって裏打ちされているわけではありません。
\item フィアットマネーは、内在的な価値がなく取引に使用することができない通貨であり、経済において本質的に価値がないものです。
\item フィアットマネーは、金や銀のような物理的な商品によって裏打ちされている通貨であり、その価値はその商品の価値に直接結びついています。
\item フィアットマネーは、中央機関の必要なしにピアツーピアの取引を可能にする分散型ネットワーク上で動作する暗号通貨の一種です。
\end{enumerate}

\bigskip


\subsection*{Question 30}
通貨の文脈において、「fiduciary」という用語は何を意味しますか？

\begin{enumerate}[label=\Alph*., leftmargin=*]
\item それは、信託取引にのみ価値がある通貨を指します。
\item それは、金に裏打ちされ、小規模な商業銀行によって発行される通貨を指します。
\item それは、それらを規制する中央機関への信頼によってのみ価値がある通貨を指します。
\item それは、特定の機関内でのみ使用される通貨を指します。
\end{enumerate}

\bigskip


\subsection*{Question 31}
信託通貨を発行する機関は何ですか？

\begin{enumerate}[label=\Alph*., leftmargin=*]
\item 政府。
\item 商業銀行。
\item 中央銀行。
\item 民間機関。
\end{enumerate}

\bigskip


\subsection*{Question 32}
Bitcoinの単位で存在できる最大数はいくつですか？

\begin{enumerate}[label=\Alph*., leftmargin=*]
\item 2100万。
\item 1億。
\item 3兆。
\item 制限はありません。
\end{enumerate}

\bigskip


\subsection*{Question 33}
Bitcoinの創造背後にある主な目的は何ですか？

\begin{enumerate}[label=\Alph*., leftmargin=*]
\item みんなを豊かにすること。
\item 金本位制を廃止すること。
\item 国家とお金を分離すること。
\item 取引をより速く、安くすること。
\end{enumerate}

\bigskip


\subsection*{Question 34}
中央機関が法定通貨の価値を下げるために使用する戦略は何ですか？

\begin{enumerate}[label=\Alph*., leftmargin=*]
\item 毎年、数パーセントずつ通貨量を増加させます。
\item 初期の通貨量と比較して、その通貨量をランダムに変動させます。
\item 毎年、数パーセントずつ通貨量を減少させます。
\item 初期の通貨量と比較して、その通貨量を一定に保ちます。
\end{enumerate}

\bigskip


\subsection*{Question 35}
通貨印刷の結果は何ですか？

\begin{enumerate}[label=\Alph*., leftmargin=*]
\item それは通貨供給の減少を引き起こします。
\item それはデフレを引き起こし、徐々に人々を豊かにします。
\item それはインフレを引き起こし、徐々に人々を貧しくします。
\item それはインフレもデフレもない安定した経済をもたらします。
\end{enumerate}

\bigskip


\subsection*{Question 36}
中央銀行デジタル通貨（CBDC）の潜在的なデメリットは何ですか？

\begin{enumerate}[label=\Alph*., leftmargin=*]
\item 金の価値の減少につながる可能性があります。
\item ビットコインの価値の減少につながる可能性があります。
\item 通貨供給の減少につながる可能性があります。
\item 個人の金融の自由を妨げ、権威主義的な乱用を容易にする可能性があります。
\end{enumerate}

\bigskip


\subsection*{Question 37}
通貨の導入において、機関がたどった歴史的な経路は何ですか？

\begin{enumerate}[label=\Alph*., leftmargin=*]
\item 指導者たちは商品のバスケットに裏打ちされた新しい通貨を導入し、その価値を安定させ続けます。
\item 指導者たちは金を中央集権化し、金と同等の価値を持つ新しい通貨を導入し、その後徐々に価値を下げていきます。
\item 指導者たちは金に裏打ちされた新しい通貨を導入し、その後徐々に価値を上げていきます。
\item 指導者たちは内在価値のない新しい通貨を導入し、その後徐々に価値を上げていきます。
\end{enumerate}

\bigskip


\subsection*{Question 38}
信託通貨に対する突然の価値下落または信頼喪失の潜在的な結果は何ですか？

\begin{enumerate}[label=\Alph*., leftmargin=*]
\item ハイパーインフレ。
\item 安定した経済。
\item デフレーション。
\item 通貨の価値の上昇。
\end{enumerate}

\bigskip


\subsection*{Question 39}
ビットコインの主要な特徴で、法定通貨と異なる点は何ですか？

\begin{enumerate}[label=\Alph*., leftmargin=*]
\item 金や銀のような物理的な商品に裏打ちされています。
\item その通貨政策は一般的な合意によって施行されます。
\item 中央銀行によって無制限に印刷されることがあります。
\item 中央銀行によって規制されています。
\end{enumerate}

\bigskip


\subsection*{Question 40}
現在の金融システムが政治的アクターにどのような潜在的影響を与えるか？

\begin{enumerate}[label=\Alph*., leftmargin=*]
\item 彼らは自身の富を確保するために現状維持を動機付けられていない。
\item 彼らは根本的な変更を行う動機付けがされておらず、システムが可能な崩壊に至るまでそのまま続けられることを許している。
\item 彼らは長期的に考えることに動機付けられている。
\item 彼らは可能な崩壊を防ぐために根本的な変更を行う動機付けがされている。
\end{enumerate}

\bigskip


\subsection*{Question 41}
ハイパーインフレーションとは何ですか？

\begin{enumerate}[label=\Alph*., leftmargin=*]
\item 通貨量のわずかな増加。
\item 通貨量のわずかな減少。
\item 通貨量の急激な減少。
\item 通貨量の急激な増加です。
\end{enumerate}

\bigskip


\subsection*{Question 42}
ハイパーインフレの結果は何ですか？

\begin{enumerate}[label=\Alph*., leftmargin=*]
\item 通貨の価値が大幅に減少します。
\item 通貨の価値が大幅に増加します。
\item 通貨の価値がわずかに減少します。
\item 通貨の価値がわずかに増加します。
\end{enumerate}

\bigskip


\subsection*{Question 43}
ハイパーインフレ現象の最初の段階は何ですか？

\begin{enumerate}[label=\Alph*., leftmargin=*]
\item 新しい通貨の出現。
\item お金の印刷の悪循環。
\item 通貨崩壊 & 価格上昇。
\item 信頼の喪失。
\end{enumerate}

\bigskip


\subsection*{Question 44}
ハイパーインフレ現象の最終段階とは何ですか？

\begin{enumerate}[label=\Alph*., leftmargin=*]
\item お金の印刷の悪循環。
\item 新しい通貨が登場します。
\item 信頼の喪失。
\item 通貨崩壊 & 価格の上昇。
\end{enumerate}

\bigskip


\subsection*{Question 45}
5年間で2%のインフレが購買力に与える影響は何ですか？

\begin{enumerate}[label=\Alph*., leftmargin=*]
\item 購買力が5%失われます。
\item 購買力が2%失われます。
\item 購買力が20%失われます。
\item 購買力が10%失われます。
\end{enumerate}

\bigskip


\subsection*{Question 46}
3年間で20%のインフレが購買力に与える影響は何ですか？

\begin{enumerate}[label=\Alph*., leftmargin=*]
\item 購買力の60%を失います。
\item 購買力の30%を失います。
\item 購買力の100%を失います。
\item 購買力の50%を失います。
\end{enumerate}

\bigskip


\subsection*{Question 47}
歴史上、1日につきインフレ率が100%に相当した事例がありますが、それはいつ、どこで起こりましたか？

\begin{enumerate}[label=\Alph*., leftmargin=*]
\item ジンバブエ、2007年。
\item 中国、1940年。
\item これは実際には起こり得ない非現実的なケースです。
\item ハンガリー、1945年。
\end{enumerate}

\bigskip


\subsection*{Question 48}
ハイパーインフレ現象の第二段階とは何ですか？

\begin{enumerate}[label=\Alph*., leftmargin=*]
\item お金の印刷の悪循環。
\item 新しい通貨の出現。
\item 通貨崩壊 & 物価上昇。
\item 信頼喪失。
\end{enumerate}

\bigskip


\subsection*{Question 49}
ハイパーインフレーションの最中、1923年10月のドイツのインフレ率はおおよそどれくらいでしたか？

\begin{enumerate}[label=\Alph*., leftmargin=*]
\item 76%
\item 110%
\item 5,200%
\item 29,500%
\end{enumerate}

\bigskip


\subsection*{Question 50}
1946年7月のハンガリーのインフレ率は何でしたか？

\begin{enumerate}[label=\Alph*., leftmargin=*]
\item 207%
\item 41,900,000,000,000,000%
\item 100%
\item 50%
\end{enumerate}

\bigskip


\subsection*{Question 51}
2008年11月のジンバブエのインフレ率は何でしたか？

\begin{enumerate}[label=\Alph*., leftmargin=*]
\item 796%
\item 79,600,000,000,000,000%
\item 100%
\item 79,600,000,000%
\end{enumerate}

\bigskip


\subsection*{Question 52}
2009年4月にジンバブエドルに何が起こったのですか？

\begin{enumerate}[label=\Alph*., leftmargin=*]
\item 政府は戦争の退役軍人に相当する4億5000万米ドルを補償することに同意しました。
\item 財務大臣はジンバブエドルの使用停止を発表し、貿易のために異なる外貨の使用を認可しました。
\item 政府は印刷機を稼働させるしかなかった。
\item ジンバブエドルは72%以上も暴落しました。
\end{enumerate}

\bigskip


\subsection*{Question 53}
次のうち、ビットコインの特徴ではないものはどれですか？

\begin{enumerate}[label=\Alph*., leftmargin=*]
\item 即時に輸送可能です。
\item 検証コストが無視できます。
\item 固定の金融政策があります。
\item 簡単に検閲されます。
\end{enumerate}

\bigskip


\subsection*{Question 54}
Bitcoinネットワーク上で取引を検証するプロセスは何と呼ばれていますか？

\begin{enumerate}[label=\Alph*., leftmargin=*]
\item 農業
\item 掘削
\item 収穫
\item マイニング
\end{enumerate}

\bigskip


\subsection*{Question 55}
マイニング報酬が半分になるイベントの名前は何ですか？

\begin{enumerate}[label=\Alph*., leftmargin=*]
\item ハービング
\item クオータリング
\item ダブリング
\item トリプリング
\end{enumerate}

\bigskip


\subsection*{Question 56}
ブロックチェーンに新しいブロックが10分ごとに追加されるメカニズムは何ですか？

\begin{enumerate}[label=\Alph*., leftmargin=*]
\item 速度調整。
\item 時間調整。
\item 難易度調整。
\item ブロック調整。
\end{enumerate}

\bigskip


\subsection*{Question 57}
人間の合理性に基づいており、ビットコインで使用される数学的概念は何ですか？

\begin{enumerate}[label=\Alph*., leftmargin=*]
\item 数論。
\item 確率論。
\item ゲーム理論。
\item 集合論。
\end{enumerate}

\bigskip


\subsection*{Question 58}
Bitcoinノード上で流通しているビットコインの量を確認するコマンドは何ですか？

\begin{enumerate}[label=\Alph*., leftmargin=*]
\item bitcoin-cli gettxoutsetinfo。
\item bitcoin-cli getblockchaininfo。
\item bitcoin-cli getnetworkinfo。
\item bitcoin-cli getwalletinfo。
\end{enumerate}

\bigskip


\subsection*{Question 59}
ビットコインはどの経済原則に沿っていますか？

\begin{enumerate}[label=\Alph*., leftmargin=*]
\item ケインズ経済学の原則。
\item マルクス経済学の原則。
\item オーストリア経済学の原則。
\item 新古典派経済学の原則。
\end{enumerate}

\bigskip


\subsection*{Question 60}
ビットコインの生成が停止されると予測されている年はいつですか？

\begin{enumerate}[label=\Alph*., leftmargin=*]
\item 2030
\item 2140
\item 2050
\item 2200
\end{enumerate}

\bigskip


\subsection*{Question 61}
難易度調整の頻度はどのくらいですか？

\begin{enumerate}[label=\Alph*., leftmargin=*]
\item 4032ブロックごと。
\item 1048ブロックごと。
\item 2016ブロックごと。
\item 512ブロックごと。
\end{enumerate}

\bigskip


\subsection*{Question 62}
マイニングの報酬はどのくらいの頻度で半減しますか？

\begin{enumerate}[label=\Alph*., leftmargin=*]
\item 540,000ブロックごと。
\item 210,000ブロックごと。
\item 120,000ブロックごと。
\item 2,100,000ブロックごと。
\end{enumerate}

\bigskip


\subsection*{Question 63}
最初の半減期後、流通していたビットコインの推定数はいくつでしたか？

\begin{enumerate}[label=\Alph*., leftmargin=*]
\item 21,000,000 BTC。
\item 5,250,000 BTC。
\item 10,500,000 BTC。
\item 15,750,000 BTC。
\end{enumerate}

\bigskip


\subsection*{Question 64}
第4回半減期後のBTC報酬はいくらですか？

\begin{enumerate}[label=\Alph*., leftmargin=*]
\item 6.25 BTC。
\item 3.125 BTC。
\item 1.5625 BTC。
\item 12.5 BTC。
\end{enumerate}

\bigskip


\subsection*{Question 65}
7回目の半減期は何年に起こりますか？

\begin{enumerate}[label=\Alph*., leftmargin=*]
\item 2036
\item 2048
\item 2028
\item 2032
\end{enumerate}

\bigskip


\subsection*{Question 66}
ビットコインウォレットの3つの主な機能は何ですか？

\begin{enumerate}[label=\Alph*., leftmargin=*]
\item ビットコインを受け取ること、ビットコインを送ること、ハッキングや盗難の試みからそれらを守ること。
\item ビットコインを受け取ることを可能にする、ビットコインを送ることを可能にする、ビットコインを売ること。
\item ビットコインを受け取ることを可能にする、ビットコインを保存する、ハッキングや盗難の試みからそれらを守ること。
\item ビットコインを保存すること、ビットコインを送ること、ハッキングや盗難の試みからそれらを守ること。
\end{enumerate}

\bigskip


\subsection*{Question 67}
Bitcoinウォレットのプライベートキーとは何ですか？

\begin{enumerate}[label=\Alph*., leftmargin=*]
\item トランザクションに署名するために使用する秘密のキーです。
\item Bitcoinアドレスを生成するために使用される公開キー。
\item 必要に応じてウォレットを破壊するための秘密の単語リスト。
\item Bitcoinトランザクションを暗号化するための秘密キー。
\end{enumerate}

\bigskip


\subsection*{Question 68}
ビットコインはどこに保存されていますか？

\begin{enumerate}[label=\Alph*., leftmargin=*]
\item ビットコインブロックチェーン内です。
\item ウォレットがインストールされているデバイス上。
\item 連邦準備制度の金庫内。
\item ビットコインウォレット内。
\end{enumerate}

\bigskip


\subsection*{Question 69}
ビットコインウォレットにおけるニーモニックフレーズとは何ですか？

\begin{enumerate}[label=\Alph*., leftmargin=*]
\item ビットコインにアクセスするための12または24の単語のリストです。
\item ビットコイン取引を暗号化するコード。
\item ビットコインウォレットにアクセスするために選択されたパスワード。
\item 必要な場合にビットコインウォレットを破壊する秘密の回復フレーズ。
\end{enumerate}

\bigskip


\subsection*{Question 70}
Bitcoinウォレットにおける公開鍵の役割は何ですか？

\begin{enumerate}[label=\Alph*., leftmargin=*]
\item Bitcoinウォレットのバックアップを取るために使用されます。
\item Bitcoin取引を暗号化するために使用されます。
\item Bitcoinウォレットにアクセスするために使用されます。
\item Bitcoinアドレスを生成し、その結果としてお金を受け取るために使用されます。
\end{enumerate}

\bigskip


\subsection*{Question 71}
Bitcoinウォレットの公開鍵を共有するリスクは何ですか？

\begin{enumerate}[label=\Alph*., leftmargin=*]
\item プライバシーの喪失。
\item Bitcoinウォレットのハッキング。
\item Bitcoinウォレットの破壊。
\item Bitcoinの喪失。
\end{enumerate}

\bigskip


\subsection*{Question 72}
以下の選択肢のうち、Bitcoinウォレットを保持しているデバイスを紛失した場合について誤っているものはどれですか？

\begin{enumerate}[label=\Alph*., leftmargin=*]
\item 公開鍵を使って、自分のBitcoinが使われていないことを確認できます。
\item ニーモニックフレーズを持っていない場合、すべてのBitcoinを失います。
\item パスワードをリセットすることで、ウォレットを再作成し、Bitcoinを使用することができます。
\item ニーモニックフレーズを使用して、ウォレットを再作成し、Bitcoinを使用することができます。
\end{enumerate}

\bigskip


\subsection*{Question 73}
ビットコインウォレットのデフォルトのセキュリティを強化するにはどうすればよいですか？

\begin{enumerate}[label=\Alph*., leftmargin=*]
\item ビットコインウォレットに第二のプライベートキーを追加することで。
\item ビットコインウォレットに第二のパブリックキーを追加することで。
\item ビットコインウォレットにパスフレーズを追加することで強化できます。
\item ビットコインウォレットに第二のニーモニックフレーズを追加することで。
\end{enumerate}

\bigskip


\subsection*{Question 74}
ビットコインウォレットの12または24の単語のリストを誰かが偶然推測する確率は？

\begin{enumerate}[label=\Alph*., leftmargin=*]
\item 極めて低い。
\item 中程度。
\item 強い。
\item 高い。
\end{enumerate}

\bigskip


\subsection*{Question 75}
ビットコインウォレットを選択する際に自分自身に問うべき中心的な質問は、自己保管の観点から何ですか？

\begin{enumerate}[label=\Alph*., leftmargin=*]
\item 自分のウォレットを他の人と共有できますか？
\item ウォレットはパスフレーズで保護されていますか？
\item ウォレットはビットコインの受け取りと送信を可能にしますか？
\item 資金の所有者はあなた自身ですか、それともお金の管理を第三者に委ねていますか？
\end{enumerate}

\bigskip


\subsection*{Question 76}
Bitcoinウォレットを使用して資金を安全に使用するために使われる技術は何ですか？

\begin{enumerate}[label=\Alph*., leftmargin=*]
\item RSA暗号。
\item 対称暗号。
\item 量子暗号。
\item 楕円曲線暗号。
\end{enumerate}

\bigskip


\subsection*{Question 77}
このコースで使用されている比喩は、送信者にあなたの資金を盗ませることなくビットコインを受け取る方法を説明するために何ですか？

\begin{enumerate}[label=\Alph*., leftmargin=*]
\item 金庫。
\item 銀行口座。
\item 家。
\item 郵便受け。
\end{enumerate}

\bigskip


\subsection*{Question 78}
ビットコインを所有している場合、主な懸念事項は何でしょうか？

\begin{enumerate}[label=\Alph*., leftmargin=*]
\item ビットコインのドル換算での価値。
\item 資金のセキュリティです。
\item ビットコインを所有することの合法性。
\item 資金を使う許可。
\end{enumerate}

\bigskip


\subsection*{Question 79}
ホットウォレットとは何ですか？

\begin{enumerate}[label=\Alph*., leftmargin=*]
\item インターネットに接続された携帯電話やコンピューター上のビットコインウォレットです。
\item ビットコインユーザーの間で人気のあるウォレット。
\item 触ると物理的に熱いウォレット。
\item 大量のビットコインを保管するために推奨されるウォレット。
\end{enumerate}

\bigskip


\subsection*{Question 80}
コールドウォレットとは何ですか？

\begin{enumerate}[label=\Alph*., leftmargin=*]
\item Bitcoinユーザーの間で不人気なウォレット。
\item 少額のビットコインを保存するために推奨されるウォレット。
\item 触れると物理的に冷たいウォレット。
\item 物理的なデバイスにあるウォレットで、キーを保存しており、インターネットに接続されていません。
\end{enumerate}

\bigskip


\subsection*{Question 81}
マルチシグウォレットとは何ですか？

\begin{enumerate}[label=\Alph*., leftmargin=*]
\item トランザクションを実行するために単一の署名が必要なウォレット。
\item トランザクションを実行するために署名が不要なウォレット。
\item 複数の署名が必要なウォレットのセットです。
\item トランザクションを実行するために物理的な署名が必要なウォレット。
\end{enumerate}

\bigskip


\subsection*{Question 82}
カストディアルサービスを使用してビットコインを保管する際の主なリスクは何ですか？

\begin{enumerate}[label=\Alph*., leftmargin=*]
\item 信頼できる第三者がいつでもあなたのビットコインの価値を下げることができます。
\item 信頼できる第三者がいつでもあなたのビットコインを複製することはできません。
\item 信頼できる第三者がいつでもあなたの資金へのアクセスを制限することができます。
\item 信頼できる第三者がいつでもあなたのビットコインの価値を上げることができます。
\end{enumerate}

\bigskip


\subsection*{Question 83}
ビットコインのセキュリティとアクセシビリティを複雑にしすぎるリスクは何ですか？

\begin{enumerate}[label=\Alph*., leftmargin=*]
\item バックアップウォレットの取り扱いを誤ると、自分自身に害を及ぼす可能性があります。
\item 物理ウォレットの取り扱いを誤ると、自分自身に害を及ぼす可能性があります。
\item マルチシグウォレットの取り扱いを誤ると、自分自身に害を及ぼす可能性があります。
\item デジタルウォレットの取り扱いを誤ると、自分自身に害を及ぼす可能性があります。
\end{enumerate}

\bigskip


\subsection*{Question 84}
モバイルウォレットの推奨される使用方法は何ですか？

\begin{enumerate}[label=\Alph*., leftmargin=*]
\item 日常の支出を保存すること。
\item 大量の金額を保存すること。
\item 長期間の貯蓄を保存すること。
\item 中期間の貯蓄を保存すること。
\end{enumerate}

\bigskip


\subsection*{Question 85}
オフライン、またはコールドの物理ウォレットの推奨される使用方法は何ですか？

\begin{enumerate}[label=\Alph*., leftmargin=*]
\item 短期間の貯蓄の保管。
\item 小額の保管。
\item 大量の保管。
\item 日常の支出の保管。
\end{enumerate}

\bigskip


\subsection*{Question 86}
マルチシグウォレットの推奨される使用方法は何ですか？

\begin{enumerate}[label=\Alph*., leftmargin=*]
\item 個人の日常費用の保管。
\item 非営利用途での中量の保管。
\item 個人用途での少量の保管。
\item 企業用途での大量の保管。
\end{enumerate}

\bigskip


\subsection*{Question 87}
パスフレーズ付きのウォレットを使用する際にバックアップする必要があるものは何ですか？

\begin{enumerate}[label=\Alph*., leftmargin=*]
\item パスフレーズのみをバックアップする必要があります。
\item 12または24の単語のリストとパスフレーズの両方をバックアップする必要があります。
\item 12または24の単語のリストのみをバックアップする必要があります。
\item 12または24の単語のリストもパスフレーズもバックアップする必要はありません。
\end{enumerate}

\bigskip


\subsection*{Question 88}
マルチシグウォレットを使用するリスクは何ですか？

\begin{enumerate}[label=\Alph*., leftmargin=*]
\item 他の当事者とプライベートキーを共有すること。
\item 他の当事者とパブリックキーを共有すること。
\item 十分な個々のプライベートキーやバックアップコンポーネントにアクセスできること。
\item 十分な個々のプライベートキーやバックアップコンポーネントにアクセスできないこと。
\end{enumerate}

\bigskip


\subsection*{Question 89}
「hodl」という用語は何を意味しますか？

\begin{enumerate}[label=\Alph*., leftmargin=*]
\item 「Hodl」は、高値で暗号通貨をすべて売却することを意味する用語です。
\item 「Hodl」は、特に市場の変動期において、投資を売却せずに保持し続けることを意味する、暗号通貨コミュニティ内の俗語です。
\item 「Hodl」は、取引のみに使用される暗号通貨の一種を指します。
\item 「Hodl」は、頻繁に購入と売却を行うデイトレーディングの戦略です。
\end{enumerate}

\bigskip


\subsection*{Question 90}
Bitcoinウォレットでプライベートキーは通常何によって表されますか？

\begin{enumerate}[label=\Alph*., leftmargin=*]
\item 12または48の単語のリスト。
\item 32の単語のリスト。
\item 64の単語のリスト。
\item 12または24の単語のリスト。
\end{enumerate}

\bigskip


\subsection*{Question 91}
私的な鍵が第三者に明らかにされた場合、どうなりますか？

\begin{enumerate}[label=\Alph*., leftmargin=*]
\item 第三者はあなたのビットコインを使うことができます。
\item 第三者はあなたの取引履歴の閲覧のみが可能です。
\item 第三者はあなたの私的な鍵を変更することができます。
\item 第三者はあなたの個人情報にアクセスできます。
\end{enumerate}

\bigskip


\subsection*{Question 92}
紙にニーモニックフレーズを保存する代わりに推奨される方法は何ですか？

\begin{enumerate}[label=\Alph*., leftmargin=*]
\item ホワイトボードに書く。
\item 電話の連絡先として保存する。
\item 金属板に刻むこと。
\item デジタルファイルに保存する。
\end{enumerate}

\bigskip


\subsection*{Question 93}
ニーモニックフレーズの単語を書き留めた後、何をすべきですか？

\begin{enumerate}[label=\Alph*., leftmargin=*]
\item 携帯電話に保存してください。
\item 自分自身にメールで送ってください。
\item 二つ目のコピーを作成し、別の場所に保管してください。
\item 信頼できる友人と共有してください。
\end{enumerate}

\bigskip


\subsection*{Question 94}
ウォレットにニーモニックフレーズがないことは何を示していますか？

\begin{enumerate}[label=\Alph*., leftmargin=*]
\item ウォレットが多くの異なる債券を保持できることを意味します。
\item ウォレットが自動的にデバイスによってバックアップされることを示唆しています。
\item ウォレットが完全に安全であり、バックアップを必要としないことを示しています。
\item それはウォレットがBitcoinの長期保管には安全ではないことを示しています。
\end{enumerate}

\bigskip


\subsection*{Question 95}
ウォレットを復元する標準的な方法は何ですか？

\begin{enumerate}[label=\Alph*., leftmargin=*]
\item ニーモニックフレーズをホワイトボードに書く。
\item 信頼できる友人とニーモニックフレーズを共有する。
\item 任意のウォレットソフトウェアにニーモニックフレーズを入力します。
\item ニーモニックフレーズをデジタルファイルに保存する。
\end{enumerate}

\bigskip


\subsection*{Question 96}
長期的にビットコインを保護する一つの方法は何ですか？

\begin{enumerate}[label=\Alph*., leftmargin=*]
\item 同僚とニーモニックフレーズを共有する。
\item 取引プラットフォーム上に保持する。
\item 鋼のような耐久性のある素材にニーモニックフレーズを刻むこと。
\item ホットウォレットに保管する。
\end{enumerate}

\bigskip


\subsection*{Question 97}
ビットコインの相続計画を作成する目的は何ですか？

\begin{enumerate}[label=\Alph*., leftmargin=*]
\item ビットコインが盗まれるのを防ぐため。
\item ビットコインの価値を下げるため。
\item 死後にビットコインが適切に管理されるようにするためです。
\item ビットコインにかかる税金を避けるため。
\end{enumerate}

\bigskip


\subsection*{Question 98}
Bitcoinの文脈で、ニーモニックフレーズではないものは何ですか？

\begin{enumerate}[label=\Alph*., leftmargin=*]
\item あなたのビットコイン資金にアクセスする秘密。
\item 特定の辞書からの24語のリスト。
\item あなたのBitcoinウォレットのパスワード。
\item 特定の辞書からの12語のリスト。
\end{enumerate}

\bigskip


\subsection*{Question 99}
ビットコインの文脈でプライバシーが重要な理由は何ですか？

\begin{enumerate}[label=\Alph*., leftmargin=*]
\item ビットコインが盗まれるのを防ぐため。
\item ビットコインの価値を下げるため。
\item 資金の追跡リスクを最小限に抑えるためです。
\item ビットコインに対する税金を支払うため。
\end{enumerate}

\bigskip


\subsection*{Question 100}
なぜビットコインを取引プラットフォームに残しておくのを避けることが推奨されるのですか？

\begin{enumerate}[label=\Alph*., leftmargin=*]
\item ハッカー攻撃の対象となりやすいからです。
\item 取引プラットフォームは規制されていない機関です。
\item 取引プラットフォームはランダムな手数料を請求します。
\item ビットコインは取引プラットフォーム上で価値が下がります。
\end{enumerate}

\bigskip


\subsection*{Question 101}
Bitcoin相続の文脈における公証人の役割は何ですか？

\begin{enumerate}[label=\Alph*., leftmargin=*]
\item 税務コンプライアンスを確保します。
\item ビットコインを金庫で安全に保管します。
\item あなたの相続計画を変更します。
\item 死後にあなたのビットコインを管理します。
\end{enumerate}

\bigskip


\subsection*{Question 102}
ビットコインウォレットの目的は何ですか？

\begin{enumerate}[label=\Alph*., leftmargin=*]
\item デバイス上にビットコインを保存するため。
\item ビットコインを取引するためにプライベートキーを保存することです。
\item デバイスのリソースを使用してビットコインをマイニングするため。
\item パブリックキーを通じてビットコインを購入するため。
\end{enumerate}

\bigskip


\subsection*{Question 103}
ビットコインにおける個人の責任の重要性は何ですか？

\begin{enumerate}[label=\Alph*., leftmargin=*]
\item ウォレットのテーマを管理する責任があるため、それが非常に重要です。
\item ビットコインを採掘する必要があるため、それが本質的に重要です。
\item ビットコインを購入することが本質的に重要です。
\item バックアップを管理する責任があるため、それが非常に重要です。
\end{enumerate}

\bigskip


\subsection*{Question 104}
Bitcoinのセキュリティの文脈で、Blockmitとは何ですか？

\begin{enumerate}[label=\Alph*., leftmargin=*]
\item 自宅でビットコインを採掘するための低コストの解決策です。
\item UTXOsを使用するために複数の署名が必要なBitcoinウォレットの一種です。
\item ニーモニックフレーズを鋼の片に刻むための低コストの解決策です。
\item 一度に複数のトランザクションに署名するためのBitcoinウォレットの一種です。
\end{enumerate}

\bigskip


\subsection*{Question 105}
KYCとは何の略ですか？

\begin{enumerate}[label=\Alph*., leftmargin=*]
\item 現金を保持せよ。
\item 顧客を知ること。
\item コードをキー入力せよ。
\item 親切に遵守せよ。
\end{enumerate}

\bigskip


\subsection*{Question 106}
なぜKYC非対応のビットコインを購入すべきですか？

\begin{enumerate}[label=\Alph*., leftmargin=*]
\item 安価で購入するため。
\item プライバシーを保護するためです。
\item 安価で売却するため。
\item すべてのビットコインユーザーの登録を保証するため。
\end{enumerate}

\bigskip


\subsection*{Question 107}
なぜ全員が同じBitcoinウォレットを使用しないのですか？

\begin{enumerate}[label=\Alph*., leftmargin=*]
\item そのようなウォレットの名前について人々が同意することは決してありません。
\item Bitcoinウォレットの使用は高価です。
\item 異なる個人は異なる使用事例を必要とします。
\item Bitcoinウォレットを使用することは違法です。
\end{enumerate}

\bigskip


\subsection*{Question 108}
ビットコインが初めて世界に紹介されたのはいつですか？

\begin{enumerate}[label=\Alph*., leftmargin=*]
\item 2010年12月12日。
\item 2009年1月3日。
\item 2009年11月22日。
\item 2008年10月31日。
\end{enumerate}

\bigskip


\subsection*{Question 109}
Satoshi Nakamotoの後でBitcoinネットワークに参加した最初の人物は誰ですか？

\begin{enumerate}[label=\Alph*., leftmargin=*]
\item Phil Champagne.
\item Laszlo Hanyecz.
\item Rogzy.
\item Hal Finney.
\end{enumerate}

\bigskip


\subsection*{Question 110}
最初に記録されたBitcoin取引は何ですか？

\begin{enumerate}[label=\Alph*., leftmargin=*]
\item Satoshi NakamotoとHal Finneyの間の10 BTCの取引です。
\item Satoshi NakamotoとHal Finneyの間の1 BTCの取引。
\item Laszlo HanyeczとHal Finneyの間の10,000 BTCの取引。
\item 2つのピザのための10,000 BTCの取引。
\end{enumerate}

\bigskip


\subsection*{Question 111}
サトシ・ナカモトはいつ姿を消しましたか？

\begin{enumerate}[label=\Alph*., leftmargin=*]
\item 2009年
\item 2008年
\item 2010年
\item 2023年
\end{enumerate}

\bigskip


\subsection*{Question 112}
ビットコインブロックチェーンの最初のブロックに含まれているメッセージは何ですか？

\begin{enumerate}[label=\Alph*., leftmargin=*]
\item 政府はNapsterのような中央集権的なネットワークの首を切るのが得意ですが、GnutellaやTorのような純粋なP2Pネットワークは自分たちの地位を保っているようです。
\item 03/jan/2009 銀行への二度目の救済措置の瀬戸際にいる財務大臣。
\item 上記のいずれでもない。
\item 武器競争の大きな戦いに勝利し、数年間の自由の新たな領域を獲得することができます。
\end{enumerate}

\bigskip


\subsection*{Question 113}
BitcoinTalkとは何ですか？

\begin{enumerate}[label=\Alph*., leftmargin=*]
\item Bitcoinのウォレット。
\item Bitcoinユーザーのための議論の場です。
\item Bitcoinのマイニングソフトウェア。
\item Bitcoinのカンファレンス。
\end{enumerate}

\bigskip


\subsection*{Question 114}
ビットコインが初めて物理的な商品の購入に使われたのはいつですか？

\begin{enumerate}[label=\Alph*., leftmargin=*]
\item Hal Finneyが100,000 BTCで車を購入したとき。
\item Laszlo Hanyeczが10,000 BTCで2枚のピザを購入したときです。
\item 他の選択肢は該当しません。
\item Satoshi Nakamotoが1,000,000,000 BTCで家を購入したとき。
\end{enumerate}

\bigskip


\subsection*{Question 115}
ビットコインネットワークが機能していたのは、その創設以来どの程度の割合ですか？

\begin{enumerate}[label=\Alph*., leftmargin=*]
\item 99.988%。
\item 80.17%。
\item 100%。
\item 90.009%。
\end{enumerate}

\bigskip


\subsection*{Question 116}
2009年以降、ビットコインのオンラインでの可用性はどの程度でしたか？

\begin{enumerate}[label=\Alph*., leftmargin=*]
\item 80.17%。
\item 90.009%。
\item 100%。
\item 99.988%。
\end{enumerate}

\bigskip


\subsection*{Question 117}
最初のBitcoin取引のブロック番号は何ですか？

\begin{enumerate}[label=\Alph*., leftmargin=*]
\item ブロック1。
\item ブロック1700。
\item ブロック170。
\item ブロック17。
\end{enumerate}

\bigskip


\subsection*{Question 118}
最初にリリースされたBitcoinソフトウェアのバージョンは何ですか？

\begin{enumerate}[label=\Alph*., leftmargin=*]
\item Bitcoin-0.0.1.
\item Bitcoin-1.0.0.
\item Bitcoin-0.1.0.
\item Bitcoin-0.1.1.
\end{enumerate}

\bigskip


\subsection*{Question 119}
サトシ・ナカモトがメールで送った最後のプライベートメッセージの日付はいつですか？

\begin{enumerate}[label=\Alph*., leftmargin=*]
\item 2010年12月12日。
\item 2011年4月23日。
\item 2008年10月31日。
\item 2009年1月8日。
\end{enumerate}

\bigskip


\subsection*{Question 120}
ビットコイン取引とは何ですか？

\begin{enumerate}[label=\Alph*., leftmargin=*]
\item 公証人を使用したビットコインの物理的な交換。
\item 電子メールアドレスを使用したビットコインの所有権の移転。
\item 物理的なアドレスを使用したビットコインの所有権の移転。
\item 暗号技術を使用したビットコインの所有権の移転です。
\end{enumerate}

\bigskip


\subsection*{Question 121}
ビットコイン取引における取引手数料の目的は何ですか？

\begin{enumerate}[label=\Alph*., leftmargin=*]
\item それは政府によって課される税金です。
\item それはウォレットプロバイダーによって請求される手数料です。
\item それはそのサービスを使用するためにビットコインネットワークによって請求される手数料です。
\item それはマイナーに次のブロックに取引を含めるためのインセンティブです。
\end{enumerate}

\bigskip


\subsection*{Question 122}
相続計画を作成する際に、含める必要がないものは何ですか？

\begin{enumerate}[label=\Alph*., leftmargin=*]
\item 資金へのアクセス方法についての説明。
\item 通常使用しているBitcoinウォレットの名前。
\item 信頼できる個人の連絡先情報。
\item 資産をリストアップした手書きの手紙。
\end{enumerate}

\bigskip


\subsection*{Question 123}
ビットコインのブロックチェーンとは何ですか？

\begin{enumerate}[label=\Alph*., leftmargin=*]
\item すべてのビットコイン取引の非公開で変更不可能な台帳です。
\item すべてのビットコイン取引の非公開で変更可能な台帳です。
\item すべてのビットコイン取引の公開されていて変更可能な台帳です。
\item すべてのビットコイン取引の公開されていて変更不可能な台帳です。
\end{enumerate}

\bigskip


\subsection*{Question 124}
ビットコイン取引における確認数が重要な理由は何ですか？

\begin{enumerate}[label=\Alph*., leftmargin=*]
\item 取引の確認数が多いほど、セキュリティが低下します。
\item 取引の確認数が多いほど、その取引はより不変性を持ちます。
\item 取引の確認数が多いほど、処理が速くなります。
\item 取引の確認数が多いほど、その取引はより価値が高まります。
\end{enumerate}

\bigskip


\subsection*{Question 125}
BitcoinネットワークにおけるBitcoinノードの役割は何ですか？

\begin{enumerate}[label=\Alph*., leftmargin=*]
\item Bitcoinネットワークを管理します。
\item 新しいブロックとその中のトランザクションを検証します。
\item Bitcoinの価格を調整します。
\item 新しいBitcoinを作成し、それらをブロックに入れます。
\end{enumerate}

\bigskip


\subsection*{Question 126}
Bitcoinノードの地理的分布の重要性は何ですか？

\begin{enumerate}[label=\Alph*., leftmargin=*]
\item ブロックチェーンの全てのコピーを破壊することが実質的に不可能になります。
\item それはBitcoin取引の速度に影響します。
\item それは一つの地域で採掘できるbitcoinsの数を決定します。
\item それは異なる地域でのBitcoinの価格を決定します。
\end{enumerate}

\bigskip


\subsection*{Question 127}
マイナーはどのようにして取引を検証し、それをブロックに含めるのですか？

\begin{enumerate}[label=\Alph*., leftmargin=*]
\item プルーフ・オブ・オーソリティのプロセスを通じて。
\item プルーフ・オブ・ワークのプロセスを通じて。
\item プルーフ・オブ・バーンのプロセスを通じて。
\item プルーフ・オブ・ステークのプロセスを通じて。
\end{enumerate}

\bigskip


\subsection*{Question 128}
Bitcoinネットワークのブロックマイニングの文脈で有効な「ハッシュ」とは何ですか？

\begin{enumerate}[label=\Alph*., leftmargin=*]
\item ブロックにアクセスするために使用されるユニークなコード。
\item ブロックに関連付けられた、256文字から成るユニークな指紋です。
\item ブロックをマイニングしたマイナーのユニークな識別子。
\item ブロックをマイニングするために必要なユニークなパスワード。
\end{enumerate}

\bigskip


\subsection*{Question 129}
Bitcoinネットワークにおける難易度調整の機能は何ですか？

\begin{enumerate}[label=\Alph*., leftmargin=*]
\item ブロックが約10分ごとにブロックチェーンに追加されるようにします。
\item Bitcoinの価格を決定します。
\item ブロックのサイズを決定します。
\item ブロックチェーンで採掘できるビットコインの数を規制します。
\end{enumerate}

\bigskip


\subsection*{Question 130}
Bitcoinプロトコル内で悪意のある行為を行うコストを高くするために、Bitcoinネットワークはどのようにしていますか？

\begin{enumerate}[label=\Alph*., leftmargin=*]
\item 法的措置の脅威を通じて。
\item 高度なセキュリティ対策の使用を通じて。
\item 厳格な規制と監視を通じて。
\item プロトコルルールと経済的インセンティブを通じて。
\end{enumerate}

\bigskip


\subsection*{Question 131}
ビットコインネットワークで、ユーザーはどのようにして自分のお金の所有権を移転しますか？

\begin{enumerate}[label=\Alph*., leftmargin=*]
\item 受取人にメールを送ることによって。
\item 私的な鍵で取引にデジタル署名をすることによって。
\item 受取人に自分のビットコインアドレスを提供することによって。
\item ビットコインを物理的に手渡すことによって。
\end{enumerate}

\bigskip


\subsection*{Question 132}
Bitcoinネットワークにおけるノードの主な機能は何ですか？

\begin{enumerate}[label=\Alph*., leftmargin=*]
\item Bitcoinブロックチェーンのコピーを保持し、トランザクションを検証し、政府に情報を伝達し、Bitcoinプロトコルのルールを変更すること。
\item 他のノードに情報を伝達することのみ。
\item Bitcoinブロックチェーンのコピーを保持することのみ。
\item Bitcoinブロックチェーンのコピーを保持し、トランザクションを検証し、他のノードに情報を伝達し、Bitcoinプロトコルのルールを施行することです。
\end{enumerate}

\bigskip


\subsection*{Question 133}
2023年9月時点で、世界中に分散しているノードのおおよその数はいくつですか？

\begin{enumerate}[label=\Alph*., leftmargin=*]
\item 45,000ノード。
\item 10,000ノード。
\item 500,000ノード。
\item 100,000ノード。
\end{enumerate}

\bigskip


\subsection*{Question 134}
現在のBitcoinブロックチェーンのサイズはどれくらいですか？

\begin{enumerate}[label=\Alph*., leftmargin=*]
\item 約2TBです。
\item 約500GBです。
\item 約100GBです。
\item 約1TBです。
\end{enumerate}

\bigskip


\subsection*{Question 135}
2023年現在、Raspberry Pi 4でBitcoinノードを運用する際の電気消費による年間コストはいくらですか？

\begin{enumerate}[label=\Alph*., leftmargin=*]
\item 年間50ユーロ強です。
\item 年間100ユーロ弱です。
\item 約200ユーロです。
\item 年間10ユーロ弱です。
\end{enumerate}

\bigskip


\subsection*{Question 136}
Bitcoinのコンテキストにおけるプルーニングされたノードとは何ですか？

\begin{enumerate}[label=\Alph*., leftmargin=*]
\item メモリ内に最後のNブロックのみを保持するノードです。
\item Bitcoinプロトコルのルールを強制するのみのノード。
\item トランザクションのみを検証するノード。
\item 他のノードに情報を伝達するのみのノード。
\end{enumerate}

\bigskip


\subsection*{Question 137}
ノードユーザーが既存の合意ルールに反する異なるルールに従うことを選択した場合、どのようなことが起こるでしょうか？

\begin{enumerate}[label=\Alph*., leftmargin=*]
\item そのノードはBitcoinネットワークから報酬を受け取ります。
\item そのノードは異なるBitcoinネットワークの一部となります。
\item そのノードはもはやBitcoinネットワークの一部ではなくなります。
\item そのノードはBitcoinネットワークから罰せられます。
\end{enumerate}

\bigskip


\subsection*{Question 138}
1ブロックが1MBで、10分ごとに生成されると考えた場合、Bitcoinノードを運用するために必要な帯域幅はどのくらいですか？

\begin{enumerate}[label=\Alph*., leftmargin=*]
\item 約10GB/月です。
\item 約1GB/月です。
\item 約5GB/月です。
\item 約20GB/月です。
\end{enumerate}

\bigskip


\subsection*{Question 139}
2017年のブロック戦争とは何でしたか？

\begin{enumerate}[label=\Alph*., leftmargin=*]
\item トランザクション容量を拡大するためにブロックサイズを増やすかどうかについての対立でした。
\item マイナーを非奨励するためにブロック報酬を減らすかどうかについての対立でした。
\item トランザクション容量を減少させるためにブロックサイズを減らすかどうかについての対立でした。
\item マイナーを奨励するためにブロック報酬を増やすかどうかについての対立でした。
\end{enumerate}

\bigskip


\subsection*{Question 140}
ビットコインノードが手頃な価格で全てのユーザーにとってアクセスしやすいままであることがなぜ重要なのか？

\begin{enumerate}[label=\Alph*., leftmargin=*]
\item それはネットワークのトランザクション容量を増加させるからです。
\item それはマイナーへのブロック報酬を増加させるからです。
\item それはネットワークのトランザクション容量を減少させるからです。
\item それはネットワークの分散化を促進するからです。
\end{enumerate}

\bigskip


\subsection*{Question 141}
ブロックが100倍重くなった場合、Bitcoinユーザーにどのような影響がありますか？

\begin{enumerate}[label=\Alph*., leftmargin=*]
\item マイナーへのブロック報酬が増加します。
\item Bitcoinノードの運用が一般の人にとってアクセスしにくくなり、プロトコルの分散化とトランザクションおよび合意ルールの不変性が損なわれます。
\item ネットワークのトランザクション容量が減少します。
\item Bitcoinノードの運用が一般の人にとってよりアクセスしやすくなり、プロトコルの分散化とトランザクションおよび合意ルールの不変性が強化されます。
\end{enumerate}

\bigskip


\subsection*{Question 142}
2017年の「ブロック戦争」の結果は何でしたか？

\begin{enumerate}[label=\Alph*., leftmargin=*]
\item ユーザーとノードはブロックサイズを増加させる提案を拒否し、SegWitと呼ばれるアップデートを活性化しました。
\item ユーザーとノードはブロックサイズを増加させる提案を拒否し、どのアップデートも活性化しませんでした。
\item ユーザーとノードはブロックサイズを増加させる提案を受け入れ、SegWitと呼ばれるアップデートを活性化しました。
\item ユーザーとノードはブロックサイズを増加させる提案を受け入れ、どのアップデートも活性化しませんでした。
\end{enumerate}

\bigskip


\subsection*{Question 143}
ライトニングネットワークとは何ですか？

\begin{enumerate}[label=\Alph*., leftmargin=*]
\item ビットコインブロックチェーンとは別のブロックチェーンネットワークです。
\item ビットコインマイナーのネットワークです。
\item ビットコインブロックチェーンに基づいた即時ビットコイン支払いネットワークです。
\item ビットコイン取引所のネットワークです。
\end{enumerate}

\bigskip


\subsection*{Question 144}
ビットコインネットワークにおけるマイナーの役割は何ですか？

\begin{enumerate}[label=\Alph*., leftmargin=*]
\item マイナーは新しいビットコインを作り出します。
\item マイナーはネットワークを保護し、トランザクションをブロックに追加します。
\item マイナーはビットコインユーザーの身元を確認します。
\item マイナーはビットコインの価格を設定します。
\end{enumerate}

\bigskip


\subsection*{Question 145}
Proof of Work (POW)とは何ですか？

\begin{enumerate}[label=\Alph*., leftmargin=*]
\item ビットコインを採掘するために使用されるアルゴリズムです。
\item 新しいビットコインを作成するプロセスです。
\item それはBitcoinプロトコルのセキュリティ合意です。
\item ビットコイン取引を検証する方法です。
\end{enumerate}

\bigskip


\subsection*{Question 146}
ハッシュレートとは何ですか？

\begin{enumerate}[label=\Alph*., leftmargin=*]
\item マイニングデバイスが1秒間に実行できるハッシュ計算の数で、プルーフ・オブ・ワークを解決するために使用されます。
\item プルーフ・オブ・ワークを解決するために、1秒間に特定の数の試行を行うことで採掘されたビットコインの総量。
\item ビットコインネットワークにおけるマイナーの総数。
\item マイニングデバイスがプルーフ・オブ・ワークを解決できる速度。
\end{enumerate}

\bigskip


\subsection*{Question 147}
Bitcoinマイニングにおける難易度調整の目的は何ですか？

\begin{enumerate}[label=\Alph*., leftmargin=*]
\item マイニングの報酬を増やすことです。
\item グローバルなマイニングゲームを参加者の数に基づいて再調整することです。
\item 流通しているビットコインの数を減らすことです。
\item Bitcoinネットワークのセキュリティを高めることです。
\end{enumerate}

\bigskip


\subsection*{Question 148}
ASICsはBitcoinマイニングにおいてどのような役割を果たしていますか？

\begin{enumerate}[label=\Alph*., leftmargin=*]
\item ASICsは、Bitcoinネットワークを保護するために使用されます。
\item ASICsは、マイニングプロセスでSHA256アルゴリズムを適用するために特別に設計されたマシンです。
\item ASICsは、検証プロセスでSHA256アルゴリズムを適用することによってBitcoinトランザクションを検証するために使用されます。
\item ASICsは、新しいbitcoinsを任意に作成するために使用されます。
\end{enumerate}

\bigskip


\subsection*{Question 149}
コインベーストランザクションとは何ですか？

\begin{enumerate}[label=\Alph*., leftmargin=*]
\item 最高の手数料が含まれるブロックの2番目のトランザクション。
\item 報酬を受け取った後にマイナーが行うトランザクション。
\item ブロックの最初のトランザクションで、マイナーの報酬が含まれます。
\item 取引所Coinbaseで行われるトランザクション。
\end{enumerate}

\bigskip


\subsection*{Question 150}
マイニングプールの目的は何ですか？

\begin{enumerate}[label=\Alph*., leftmargin=*]
\item マイニングプールは、全体のBitcoinネットワークのハッシュレートを増加させるために使用されます。
\item マイニングプールは、マイニングの報酬を増加させるために使用されます。
\item マイニングプールは、マイナーが計算リソースをプールし、より小さいがより定期的な報酬を受け取ることを可能にします。
\item マイニングプールは、マイニングの難易度を下げるために使用されます。
\end{enumerate}

\bigskip


\subsection*{Question 151}
ビザンチン将軍問題とは何ですか？

\begin{enumerate}[label=\Alph*., leftmargin=*]
\item 中央集権的ネットワークでの合意形成を維持する問題です。
\item それは、様々なアクターを信頼できないシステムで情報を調整する問題です。
\item 分散ネットワークでのトランザクションの検証問題です。
\item デジタル通貨での二重支払いを防ぐ問題です。
\end{enumerate}

\bigskip


\subsection*{Question 152}
BitcoinネットワークにおけるProof of Workの目的は何ですか？

\begin{enumerate}[label=\Alph*., leftmargin=*]
\item Proof of Workは、第三者を信用して新しいビットコインを作成するために使用されます。
\item Proof of Workは、第三者を信用することなく情報の真実性を保証します。
\item Proof of Workは、第三者を信用してビットコイン取引を検証するために使用されます。
\item Proof of Workは、第三者を信用することなくビットコインプロトコルを変更するために使用されます。
\end{enumerate}

\bigskip


\subsection*{Question 153}
ビットコインマイニングにおけるエネルギー消費の意義は何ですか？

\begin{enumerate}[label=\Alph*., leftmargin=*]
\item エネルギーコストはマイニングに使用されるASICマシンを動かすために必要であり、情報の検証には高コストがかかります。
\item 新しいビットコインを作り出すためにエネルギーコストは必要ありません。
\item エネルギー消費は情報を生産するために必要ですが、情報の検証コストは無視できます。この非対称性がネットワークのセキュリティを保証します。
\item エネルギー消費はビットコインネットワークを維持するために必要であり、情報の検証には可変コストがかかります。
\end{enumerate}

\bigskip


\subsection*{Question 154}
ビットコインネットワークで51%攻撃が起こった場合、何が起こりますか？

\begin{enumerate}[label=\Alph*., leftmargin=*]
\item エージェントが流通しているビットコインの半分以上を支配下に置く。
\item エージェントがビットコインネットワークのマイニングプールの半分以上を支配下に置く。
\item エージェントがハッシュレートの半分以上を持ち、ブロックチェーンを変更しようと試みることができます。
\item エージェントがビットコインネットワークのノードの半分以上を支配下に置く。
\end{enumerate}

\bigskip


\subsection*{Question 155}
ビットコインマイニングの環境コストに関する主な懸念は何ですか？

\begin{enumerate}[label=\Alph*., leftmargin=*]
\item 電気のコスト。
\item エネルギー消費。
\item ASICのリサイクル。
\item ASICの生産。
\end{enumerate}

\bigskip


\subsection*{Question 156}
Bitcoinプロトコルにおけるマイナーの役割は何ですか？

\begin{enumerate}[label=\Alph*., leftmargin=*]
\item 現在および過去のネットワークを保護します。
\item Bitcoinの価格を調整します。
\item Bitcoinの供給を制御します。
\item ユーザー間の取引を管理します。
\end{enumerate}

\bigskip


\subsection*{Question 157}
Bitcoinプロトコルにおいて、2つのブロックが生成される平均時間はどのくらいですか？

\begin{enumerate}[label=\Alph*., leftmargin=*]
\item 20分。
\item 10分。
\item 5分。
\item 15分。
\end{enumerate}

\bigskip


\subsection*{Question 158}
金融的抑圧や独裁政権下に住む個人にとって、ビットコインの主な利点は何ですか？

\begin{enumerate}[label=\Alph*., leftmargin=*]
\item 投資による高いリターンを得ること。
\item 固定された価値を持つ安定した通貨を使用すること。
\item 検閲や銀行の制限から逃れるチャンスを持つこと。
\item 匿名で取引を行うことを可能にすること。
\end{enumerate}

\bigskip


\subsection*{Question 159}
Bitcoinプロトコルは、2つのブロックの作成間の平均時間をどのように一定に保っていますか？

\begin{enumerate}[label=\Alph*., leftmargin=*]
\item 有効なハッシュを見つける難易度は、2016ブロックごとに調整されます。
\item 有効なハッシュを見つけた報酬が増加します。
\item ブロックのサイズは10メガバイトに調整されます。
\item マイナーの数は固定数に調整されます。
\end{enumerate}

\bigskip


\subsection*{Question 160}
まだグリッドに接続されていない発電所に対して、鉱夫は何をすることができますか？

\begin{enumerate}[label=\Alph*., leftmargin=*]
\item 彼らは発電所をグリッドに接続することができます。
\item 彼らは発電所に電力を提供することができます。
\item 彼らは電気の価格を規制することができます。
\item 彼らは最後の手段としての買い手として行動することができます。
\end{enumerate}

\bigskip


\subsection*{Question 161}
現在の金融システムが環境に対して抱える主な問題は何ですか？

\begin{enumerate}[label=\Alph*., leftmargin=*]
\item 商品の生産を規制していない。
\item 再生可能エネルギーへの投資がない。
\item 過剰消費と借金を助長していることです。
\item 化石燃料の使用を促進している。
\end{enumerate}

\bigskip


\subsection*{Question 162}
ビットコインはどのようにして生態系の移行を助ける可能性がありますか？

\begin{enumerate}[label=\Alph*., leftmargin=*]
\item 資源の消費を規制します。
\item グリーンエネルギーの使用を促進します。
\item 商品の生産を減らします。
\item 温室効果ガスの排出のみを制限します。
\end{enumerate}

\bigskip


\subsection*{Question 163}
ビットコインプロトコルは、次のブロックの有効なハッシュを見つけたマイナーにどのように報酬を与えますか？

\begin{enumerate}[label=\Alph*., leftmargin=*]
\item 合意ルールによって定義された報酬と、有効なブロックに含まれるすべてのトランザクションからの手数料で報酬を与えます。
\item 有効なブロックに含まれるすべてのトランザクションからの手数料で報酬を与えます。
\item 総ビットコイン供給量の割合と、前のブロックに含まれるすべてのトランザクションからの手数料で報酬を与えます。
\item 合意ルールによって定義されていない一定量のビットコインで報酬を与えます。
\end{enumerate}

\bigskip


\subsection*{Question 164}
Bitcoinマイニングにおける専用SHA256マシン、またはASICの主な利点は何ですか？

\begin{enumerate}[label=\Alph*., leftmargin=*]
\item トランザクションの速度を上げます。
\item ジュールあたりのハッシュレートを増加させます。
\item エネルギー消費を減らします。
\item ネットワークのセキュリティを強化します。
\end{enumerate}

\bigskip


\subsection*{Question 165}
ビットコインプロトコルは、どのようにして検閲耐性と回復力を保証しているのですか？

\begin{enumerate}[label=\Alph*., leftmargin=*]
\item トランザクションを検証するために複雑なアルゴリズムを使用しています。
\item 高レベルの暗号化を使用しています。
\item ヨーロッパに分散されたノードのローカルネットワークを使用しています。
\item プロトコルの各コンポーネントは、地球上の地理的に分散された場所に配置されています。
\end{enumerate}

\bigskip


\subsection*{Question 166}
金融システムに関して、ケインズ経済学とオーストリア経済学の主な違いは何ですか？

\begin{enumerate}[label=\Alph*., leftmargin=*]
\item ケインズ経済学は、状況や資源の時間的および動的な側面を考慮に入れていません。
\item ケインズ経済学はインフレを奨励します。
\item ケインズ経済学は通貨の過剰消費を促進します。
\item ケインズ経済学は状況や資源の希少性を認識しています。
\end{enumerate}

\bigskip


\subsection*{Question 167}
ビットコインが特定のグループに魅力的である理由の一つは何ですか？

\begin{enumerate}[label=\Alph*., leftmargin=*]
\item 中央権力によって制御され、資本をユーザーに自由に配布される。
\item 匿名性がないため、すべてのユーザーがよく知られ、信頼されている。
\item 信頼できる電子マネー、検閲への抵抗力、透明で不変の金融政策などの解決策を提供できる能力です。
\item 偽造が容易で、クリック一つでより多くのお金を得る解決策を作り出せる。
\end{enumerate}

\bigskip


\subsection*{Question 168}
マイナーにとっての典型的なコストは何ですか？

\begin{enumerate}[label=\Alph*., leftmargin=*]
\item 電力消費とハードウェアの購入。
\item ハードウェアの購入のみ。
\item 電力消費とトランザクション手数料。
\item トランザクション手数料。
\end{enumerate}

\bigskip


\subsection*{Question 169}
ビットコインの価格変動を軽減する一つの方法は何ですか？

\begin{enumerate}[label=\Alph*., leftmargin=*]
\item ファイナンシャルヘッジ（ステーブルコイン）。
\item すべてのお金をビットコインに投資する。
\item 市場のトレンドを無視する。
\item 価格が下がったときにさらにビットコインを購入する。
\end{enumerate}

\bigskip


\subsection*{Question 170}
Bitcoinプロトコルのユニークな特徴は何ですか？

\begin{enumerate}[label=\Alph*., leftmargin=*]
\item それは簡単にハッキングされます。
\item それはビジネス時間中のみ稼働します。
\item それは世界規模で、1日24時間、週7日間稼働しています。
\item それは中央権力によって制御されています。
\end{enumerate}

\bigskip


\subsection*{Question 171}
2017年に暗号通貨界で起きた重要な出来事は何ですか？

\begin{enumerate}[label=\Alph*., leftmargin=*]
\item ビットコインの禁止。
\item 何千もの初期コインオファリング（ICO）の開始。
\item すべての暗号通貨の消失。
\item ビットコインの創造。
\end{enumerate}

\bigskip


\subsection*{Question 172}
ビットコインの全体的な傾向は、限定された供給量と半減期プロセスを考慮すると、どのようなものですか？

\begin{enumerate}[label=\Alph*., leftmargin=*]
\item 一般的に下向きの傾向。
\item 傾向がまったくない。
\item 一般的に上向きの傾向。
\item 識別可能なパターンなしに激しく変動する傾向。
\end{enumerate}

\bigskip


\subsection*{Question 173}
ビットコインが多くの投機的なバブルを経験する理由の一つは何ですか？

\begin{enumerate}[label=\Alph*., leftmargin=*]
\item ビットコインは経済サイクルの影響を受けないからです。
\item ビットコインは世界中どこでも使用される通貨だからです。
\item ビットコインは実際にはどこにも使用されていないからです。
\item ビットコインは消えてしまうバブルだからです。
\end{enumerate}

\bigskip


\subsection*{Question 174}
金融当局がビットコインを規制することを難しくしている一因は何ですか？

\begin{enumerate}[label=\Alph*., leftmargin=*]
\item 中央機関によって管理されています。
\item 世界のどこでも使用されていません。
\item それは世界規模で、1日24時間、週7日間稼働しています。
\item 偽造が容易です。
\end{enumerate}

\bigskip


\subsection*{Question 175}
ビットコインが従来の市場にさらに統合された一つの方法は何ですか？

\begin{enumerate}[label=\Alph*., leftmargin=*]
\item すべての暗号通貨の消失。
\item ビットコインの創造。
\item ビットコインの禁止。
\item ビットコインETFの登場。
\end{enumerate}

\bigskip


\subsection*{Question 176}
次のうち、Silk Roadのようなダークウェブマーケットプレイスでビットコインの使用が拡大した要因はどれですか？

\begin{enumerate}[label=\Alph*., leftmargin=*]
\item 追跡可能で制御可能な性質。
\item 匿名性の欠如。
\item 制御不可能で匿名性がある性質。
\item 中央集権的な制御。
\end{enumerate}

\bigskip


\subsection*{Question 177}
ビットコインの価格が数日で10%、20%、あるいは50%も下落する一つの理由は何ですか？

\begin{enumerate}[label=\Alph*., leftmargin=*]
\item 強気市場と弱気市場の相場変動。
\item 金の価格。
\item 天気。
\item インターネットを使用する人の数。
\end{enumerate}

\bigskip


\subsection*{Question 178}
ビットコインが消えるバブルではない理由の一つは何ですか？

\begin{enumerate}[label=\Alph*., leftmargin=*]
\item 実際にはどこでも使用されていないからです。
\item 偽造が容易だからです。
\item それは世界中どこでも実際に使用されている通貨だからです。
\item 中央権力によって制御されているからです。
\end{enumerate}

\bigskip


\subsection*{Question 179}
フィアット通貨と比較して、ビットコインはどのように捉えられていますか？

\begin{enumerate}[label=\Alph*., leftmargin=*]
\item 第三次経済として。
\item 産業経済として。
\item 並行経済として。
\item 代替経済として。
\end{enumerate}

\bigskip


\subsection*{Question 180}
ビットコインを購入せずに入手する方法の一つは何ですか？

\begin{enumerate}[label=\Alph*., leftmargin=*]
\item 提供する製品やサービスの支払い手段として受け入れる。
\item 友人から借りて返さない。
\item 宝くじで当てる。
\item 通りで見つける。
\end{enumerate}

\bigskip


\subsection*{Question 181}
マーチャントとしてビットコインを受け入れる一つの利点は何ですか？

\begin{enumerate}[label=\Alph*., leftmargin=*]
\item デフレーションからの保護。
\item 検閲への抵抗。
\item 取引手数料の増加。
\item 財政制限。
\end{enumerate}

\bigskip


\subsection*{Question 182}
BTCMapとは何ですか？

\begin{enumerate}[label=\Alph*., leftmargin=*]
\item 世界中のすべてのBitcoin取引所の地図です。
\item Bitcoinを受け入れるすべての加盟店をリストアップしたプラットフォームです。
\item 世界中のすべてのBitcoinユーザーの地図です。
\item 世界中のすべてのBitcoin鉱山の地図です。
\end{enumerate}

\bigskip


\subsection*{Question 183}
ビットコインはどのようにしてビザンチン将軍問題を解決していますか？

\begin{enumerate}[label=\Alph*., leftmargin=*]
\item プルーフ・オブ・ワークに基づくブロックチェーンと、変動する難易度を使用することで。
\item デジタル署名の使用を通じて。
\item 公開台帳の使用を通じて。
\item プルーフ・オブ・ステークアルゴリズムの使用を通じて。
\end{enumerate}

\bigskip


\subsection*{Question 184}
AMFとは何ですか？

\begin{enumerate}[label=\Alph*., leftmargin=*]
\item 金融市場庁。
\item 自動マネーフロー。
\item アメリカ通貨基金。
\item 鉱業施設協会。
\end{enumerate}

\bigskip


\subsection*{Question 185}
ビジネスでビットコインを受け入れるソリューションを選択する際に考慮すべき1つの要因は何ですか？

\begin{enumerate}[label=\Alph*., leftmargin=*]
\item 従業員の数。
\item 本社の位置。
\item 予想される取引量。
\item ロゴの色。
\end{enumerate}

\bigskip


\subsection*{Question 186}
OpenNodeとは何ですか？

\begin{enumerate}[label=\Alph*., leftmargin=*]
\item ビットコイン交換プラットフォームです。
\item ビジネスがビットコインを受け入れるためのシンプルなオンラインソリューションです。
\item オープンソースのビットコインマイニングソフトウェアです。
\item 多くのユーザーによって同時に使用されるビットコインノードの一種です。
\end{enumerate}

\bigskip


\subsection*{Question 187}
ビットコイン経済をよりアクセスしやすく、便利にするために貢献している取り組みの一つは何ですか？

\begin{enumerate}[label=\Alph*., leftmargin=*]
\item ビットコインマイニング。
\item ビットコイン取引所。
\item BTCMap。
\item ビットコインロト。
\end{enumerate}

\bigskip


\subsection*{Question 188}
大規模な組織や情熱的なBitcoin愛好家がBitcoinを受け入れるためのツールは何ですか？

\begin{enumerate}[label=\Alph*., leftmargin=*]
\item Bitcoin Exchange.
\item Bitcoin Wallet.
\item BTCpay Server.
\item BTCpay Mining.
\end{enumerate}

\bigskip


\subsection*{Question 189}
次のうち、日常生活でビットコインを使用したい人を助ける行動はどれですか？

\begin{enumerate}[label=\Alph*., leftmargin=*]
\item ビットコインの価格を上げる。
\item ビットコインの使用を制限する。
\item ビットコインの数を減らす。
\item ビットコインを受け入れるすべての商人のリストを持つ。
\end{enumerate}

\bigskip


\subsection*{Question 190}
Swiss Bitcoin Payとは何ですか？

\begin{enumerate}[label=\Alph*., leftmargin=*]
\item ビットコインを受け入れるスイスの銀行です。
\item スイスで活動する鉱業会社です。
\item アマチュアの商人がビットコインを受け入れるためのソリューションです。
\item スイス発のビットコイン交換所です。
\end{enumerate}

\bigskip


\subsection*{Question 191}
ビットコインを購入する際に関連するリスクの一つは何ですか？

\begin{enumerate}[label=\Alph*., leftmargin=*]
\item 物理的に失われることがあります。
\item その価格が0に下がる可能性があります。
\item 簡単に盗まれることがあります。
\item 複製されることがあります。
\end{enumerate}

\bigskip


\subsection*{Question 192}
Bitcoinを購入する前に持っておくべきものの一つは何ですか？

\begin{enumerate}[label=\Alph*., leftmargin=*]
\item クレジットカード。
\item 物理的な金庫。
\item 銀行口座。
\item 安全なウォレット。
\end{enumerate}

\bigskip


\subsection*{Question 193}
ドルコスト平均法とは何ですか？

\begin{enumerate}[label=\Alph*., leftmargin=*]
\item 価格が下がったときにのみビットコインを購入すること。
\item 定期的な間隔で少額のビットコインを購入すること。
\item 価格が上がったときにのみビットコインを購入すること。
\item 一度に大量のビットコインを購入すること。
\end{enumerate}

\bigskip


\subsection*{Question 194}
自発的な購入とは何ですか？

\begin{enumerate}[label=\Alph*., leftmargin=*]
\item 価格が上がったときにのみビットコインを購入すること。
\item 即座にビットコインを購入して、その影響を受けること。
\item 価格が下がったときにのみビットコインを購入すること。
\item 定期的な間隔でビットコインを購入すること。
\end{enumerate}

\bigskip


\subsection*{Question 195}
KYC規制の目的は何ですか？

\begin{enumerate}[label=\Alph*., leftmargin=*]
\item テロ資金の調達、脱税、およびマネーロンダリングと戦うこと。
\item ビットコインの価格を上げて、ユーザーにより多くの税金を払わせるため。
\item ビットコインの供給を制御するため。
\item 人々がビットコインを購入するのを防ぐため。
\end{enumerate}

\bigskip


\subsection*{Question 196}
KYCなしでビットコインを取得する方法の一つは何ですか？

\begin{enumerate}[label=\Alph*., leftmargin=*]
\item ブローカープラットフォーム。
\item DCAプラットフォーム。
\item ビットコインATM。
\item 銀行振込。
\end{enumerate}

\bigskip


\subsection*{Question 197}
取引プラットフォームでビットコインを購入した後、どのような行動を取るべきですか？

\begin{enumerate}[label=\Alph*., leftmargin=*]
\item プラットフォーム上にビットコインを残しておく。
\item すぐにビットコインを売却する。
\item ビットコインを引き出して、自分のウォレットに送金してください。
\item ビットコインを別の取引プラットフォームに転送する。
\end{enumerate}

\bigskip


\subsection*{Question 198}
DCAプラットフォームを選択する際に考慮すべきことの1つは何ですか？

\begin{enumerate}[label=\Alph*., leftmargin=*]
\item プラットフォームのロゴの色。
\item プラットフォーム上のユーザー数。
\item あなたの国での利用可能性。
\item プラットフォームの年齢。
\end{enumerate}

\bigskip


\subsection*{Question 199}
時々あなたのウォレット内のUTXOを統合する理由の一つは何ですか？

\begin{enumerate}[label=\Alph*., leftmargin=*]
\item ビットコインの価値を高めるため。
\item ビットコインが常に利用可能であることを確認するため。
\item 取引中の不必要な手数料を避けるためです。
\item ビットコインが盗まれるのを防ぐため。
\end{enumerate}

\bigskip


\subsection*{Question 200}
FOMOとFUDとは何ですか？

\begin{enumerate}[label=\Alph*., leftmargin=*]
\item ビットコインを購入する戦略です。
\item ビットコインに関連する規制です。
\item ビットコインウォレットの種類です。
\item 衝動的で潜在的に有害な意思決定につながる可能性のある感情です。
\end{enumerate}

\bigskip


\subsection*{Question 201}
取引所プラットフォームにビットコインを残しておくことの潜在的なリスクは何ですか？

\begin{enumerate}[label=\Alph*., leftmargin=*]
\item ビットコインの価格が他の暗号通貨によって上回られるリスク。
\item ハッキングと資金ブロックのリスクです。
\item プラットフォームがビットコインの供給を制御するリスク。
\item プラットフォームがビットコインの価格を上げるリスク。
\end{enumerate}

\bigskip


\subsection*{Question 202}
あなたの国の規制に従わないでビットコインを購入する際の潜在的な結果の一つは何ですか？

\begin{enumerate}[label=\Alph*., leftmargin=*]
\item 他の人がビットコインを購入するのを防ぐ可能性があります。
\item ビットコインの供給をコントロールする可能性があります。
\item ビットコインの価格を上げる可能性があります。
\item 自分自身を危険にさらす可能性があります。
\end{enumerate}

\bigskip


\subsection*{Question 203}
ビットコインの採用曲線は何に似ていますか？

\begin{enumerate}[label=\Alph*., leftmargin=*]
\item Sカーブ。
\item Uカーブ。
\item Vカーブ。
\item Wカーブ。
\end{enumerate}

\bigskip


\subsection*{Question 204}
どの国がビットコインを法定通貨として宣言しましたか？

\begin{enumerate}[label=\Alph*., leftmargin=*]
\item エルサルバドル。
\item アメリカ合衆国。
\item 中国。
\item インド。
\end{enumerate}

\bigskip


\subsection*{Question 205}
ビットコインは企業、大学、規制機関、個人にどのような行動を強いるのか？

\begin{enumerate}[label=\Alph*., leftmargin=*]
\item それを考慮に入れること。
\item それを犯罪化すること。
\item それを無視すること。
\item それを禁止すること。
\end{enumerate}

\bigskip


\subsection*{Question 206}
このコースはビットコインをどのように説明していますか？

\begin{enumerate}[label=\Alph*., leftmargin=*]
\item 多分野にわたる科目。
\item 単一分野の科目。
\item 物理的な科目。
\item 単純な科目。
\end{enumerate}

\bigskip


\subsection*{Question 207}
このコースによると、ビットコイン革命の主な目標は何ですか？

\begin{enumerate}[label=\Alph*., leftmargin=*]
\item 現金が時代遅れになる世界を創造すること。
\item 政府に権力がなく、社会が混沌とする世界を創造すること。
\item プライバシーが人権として認められず、銀行によってお金が操作される世界を創造すること。
\item 人間の主権が権利として認められ、プライバシーがデフォルトで尊重され、お金が操作されない世界を創造すること。
\end{enumerate}

\bigskip


\subsection*{Question 208}
1999年にミルトン・フリードマンは何を予測しましたか？

\begin{enumerate}[label=\Alph*., leftmargin=*]
\item 政府によるお金の管理の終焉。
\item 信頼性の高い電子マネーの開発。
\item ビットコインの台頭。
\item 現金の消滅。
\end{enumerate}

\bigskip


\subsection*{Question 209}
Bitcoinが一般的な上昇傾向にあるとはどういう意味ですか？

\begin{enumerate}[label=\Alph*., leftmargin=*]
\item Bitcoinの一般的な上昇傾向とは、価格が無限に上昇し続けることが保証されており、どんな変動もないことを意味します。
\item Bitcoinの一般的な上昇傾向とは、それが少数の投資家の間でのみ人気があり、全体の市場には影響がないことを意味します。
\item Bitcoinの一般的な上昇傾向とは、Bitcoinが価値を失い、市場の魅力を失っていることを意味します。
\item Bitcoinの一般的な上昇傾向とは、全体として、Bitcoinの価格が増加していることを意味します。
\end{enumerate}

\bigskip


\subsection*{Question 210}
エルサルバドルにおけるビットコインの地位は何ですか？

\begin{enumerate}[label=\Alph*., leftmargin=*]
\item 犯罪化されている。
\item 規制されていない。
\item 不明。
\item 法定通貨。
\end{enumerate}

\bigskip


\subsection*{Question 211}
ビットコインはどの分野で疑問を提起しますか？

\begin{enumerate}[label=\Alph*., leftmargin=*]
\item コンピュータ科学と経済学のみ
\item 暗号学とゲーム理論のみ。
\item 哲学とエネルギーのみ。
\item 暗号学、ゲーム理論、経済学と金融政策、コンピュータ科学、哲学、エネルギー、法律、および規制。
\end{enumerate}

\bigskip


\subsection*{Question 212}
このコースで説明されているビットコインの普及の性質は何ですか？

\begin{enumerate}[label=\Alph*., leftmargin=*]
\item 緩やかで安定している。
\item 急速で制御不可能。
\item 断続的で予測不可能。
\item バイラル。
\end{enumerate}

\bigskip


\subsection*{Question 213}
このコースはビットコインの未来について何を示唆していますか？

\begin{enumerate}[label=\Alph*., leftmargin=*]
\item より良い技術に置き換えられるでしょう。
\item その価値を失うでしょう。
\item それは消え去ることはないでしょう。
\item 世界中で禁止されるでしょう。
\end{enumerate}

\bigskip


\subsection*{Question 214}
このコースは、ビットコインについて学ぶペースについて何を示唆していますか？

\begin{enumerate}[label=\Alph*., leftmargin=*]
\item 一度にすべてを理解するのは複雑なので、時間をかけてください。
\item 一度にすべてを簡単に学ぶことができます。
\item ビットコインについてすべてを学ぶ必要はありません。
\item ビットコインについてすべてを学ぶことは不可能です。
\end{enumerate}

\bigskip


\subsection*{Question 215}
Lightning Networkとは何ですか？

\begin{enumerate}[label=\Alph*., leftmargin=*]
\item Bitcoinプロトコルを使用して、稲妻のように速い取引を可能にする支払いネットワークです。
\item 新しいタイプの暗号通貨です。
\item 将来的にBitcoinを置き換える新しいブロックチェーン技術です。
\item Bitcoinを採掘する方法です。
\end{enumerate}

\bigskip


\subsection*{Question 216}
Bitcoinのスケーラビリティ問題とは何ですか？

\begin{enumerate}[label=\Alph*., leftmargin=*]
\item 時間が経つにつれて新しいBitcoinを採掘する難しさが増すこと。
\item Bitcoinに対する規制の欠如が取引を違法にすること。
\item 採用されるにつれて、秒間の取引数を増やし続けることができる金融システムを実装するという課題です。
\item ハービングの後にBitcoinの価格の変動性が増すこと。
\end{enumerate}

\bigskip


\subsection*{Question 217}
ブロックチェーンのトリレンマとは何ですか？

\begin{enumerate}[label=\Alph*., leftmargin=*]
\item ブロックチェーンに基づくプロトコルは、分散化、セキュリティ、スケーラビリティのいずれの基準も満たすことができません。
\item ブロックチェーンに基づくプロトコルは、分散化、セキュリティ、スケーラビリティの基準のうち2つしか満たすことができません。
\item ブロックチェーンに基づくプロトコルは、分散化、セキュリティ、スケーラビリティの全ての基準を満たすことができます。
\item ブロックチェーンに基づくプロトコルは、分散化、セキュリティ、スケーラビリティの基準のうち1つしか満たすことができません。
\end{enumerate}

\bigskip


\subsection*{Question 218}
Lightning Networkは何を可能にしますか？

\begin{enumerate}[label=\Alph*., leftmargin=*]
\item 即時で低コストのBitcoin取引。
\item 即時で高コストのBitcoin取引。
\item 遅延した低コストのBitcoin取引。
\item 遅延した高コストのBitcoin取引。
\end{enumerate}

\bigskip


\subsection*{Question 219}
Lightning Networkはいつ最初に検証され、実装されましたか？

\begin{enumerate}[label=\Alph*., leftmargin=*]
\item 2017
\item 2016
\item 2018
\item 2015
\end{enumerate}

\bigskip


\subsection*{Question 220}
Lightning Networkはどのようなネットワークを作り出しますか？

\begin{enumerate}[label=\Alph*., leftmargin=*]
\item 支払いチャネルのネットワーク。
\item ブロックチェーンノードのネットワーク。
\item 採掘ノードのネットワーク。
\item 仮想通貨取引所のネットワーク。
\end{enumerate}

\bigskip


\subsection*{Question 221}
従来の送金サービス（Western Union、中央銀行、Visa、Mastercardなど）がLightning Networkを採用しなかった場合、どのようなことが起こり得るでしょうか？

\begin{enumerate}[label=\Alph*., leftmargin=*]
\item 彼らは消滅する可能性があります。
\item より人気が出る可能性があります。
\item より分散化する可能性があります。
\item より安全になる可能性があります。
\end{enumerate}

\bigskip


\subsection*{Question 222}
Lightning Network上の取引手数料は非常に低く、多くの場合...以下ですか？

\begin{enumerate}[label=\Alph*., leftmargin=*]
\item 3%未満。
\item 1%未満。
\item 0.5%未満。
\item 2%未満。
\end{enumerate}

\bigskip


\subsection*{Question 223}
Lightning Network上のトランザクションはどのようにして保護されていますか？

\begin{enumerate}[label=\Alph*., leftmargin=*]
\item 公開鍵のみを使用して。
\item 暗号技術を通じて、および間接的にはBitcoin上のマイナーによって消費されるエネルギーを通じて。
\item 私的鍵のみを使用して。
\item 私的鍵と公開鍵の両方のみを使用して。
\end{enumerate}

\bigskip


\subsection*{Question 224}
Lightning Networkは何を可能にしますか？

\begin{enumerate}[label=\Alph*., leftmargin=*]
\item 直接のチャネル接続を持たない個人間での取引。
\item 大規模な取引に限定された直接チャネル接続を確立すること。
\item 中央集権型の取引所とのみ直接チャネル接続を作成すること。
\item 任意の取引に複数の確認が必要なチャネル接続を作成すること。
\end{enumerate}

\bigskip


\subsection*{Question 225}
Bitcoinプロトコルにおいて、2つのブロックが生成される平均時間間隔は何分ですか？

\begin{enumerate}[label=\Alph*., leftmargin=*]
\item 20分。
\item 5分。
\item 15分。
\item 10分。
\end{enumerate}

\bigskip


\subsection*{Question 226}
Bitcoinプロトコルにおけるブロックサイズはどのくらいですか？

\begin{enumerate}[label=\Alph*., leftmargin=*]
\item 2MB。
\item 3MB。
\item 4MB。
\item 1MB。
\end{enumerate}

\bigskip


\subsection*{Question 227}
Lightning Networkの主な目的は何ですか？

\begin{enumerate}[label=\Alph*., leftmargin=*]
\item Bitcoinの価値を高めること。
\item マイクロトランザクションを容易にすること。
\item 新しい暗号通貨を作成すること。
\item Bitcoinプロトコルを置き換えること。
\end{enumerate}

\bigskip


\subsection*{Question 228}
Lightning Networkの日常生活における潜在的な応用の一つは何ですか？

\begin{enumerate}[label=\Alph*., leftmargin=*]
\item コマースにおける即時請求。
\item 新しいビットコインの採掘。
\item ビットコインの価値を高める。
\item 新しい暗号通貨の作成。
\end{enumerate}

\bigskip


\subsection*{Question 229}
Lightning Networkで使用されるBitcoinの単位は何ですか？

\begin{enumerate}[label=\Alph*., leftmargin=*]
\item Block。
\item Sats。
\item Ether。
\item Litecoin。
\end{enumerate}

\bigskip


\subsection*{Question 230}
ビデオゲーム業界における「skin in the game」の概念とは何ですか？

\begin{enumerate}[label=\Alph*., leftmargin=*]
\item ゲームの結果に対する財務的な利害関係を持つこと。
\item ゲームプレイを向上させるユニークなスキンでキャラクターをカスタマイズすること。
\item 実際のお金でゲームのスキンを売ること。
\item 実際のお金でゲーム内アイテムを購入すること。
\end{enumerate}

\bigskip


\subsection*{Question 231}
Lightning Networkは、ビデオゲーム業界の既存のビジネスモデルをどのように変革する可能性がありますか？

\begin{enumerate}[label=\Alph*., leftmargin=*]
\item プレイヤーがゲームをプレイしながらBitcoinをマイニングできるようになります。
\item プレイヤーがゲームをプレイする際に、非常に少額のお金を賭けることができます。
\item プレイヤーがBitcoinでゲーム内アイテムを購入できるようになります。
\item ゲーム開発者が自分たちのゲームに対してBitcoinでの支払いを受け入れられるようになります。
\end{enumerate}

\bigskip


\subsection*{Question 232}
Lightning Networkの文脈での「マネーストリーミング」とは何ですか？

\begin{enumerate}[label=\Alph*., leftmargin=*]
\item ビットコインでオーディオやビデオコンテンツをストリーミングすること。
\item 継続的でリアルタイムのマイクロペイメントを行うこと。
\item ATMを使用してビットコインを現金に変換すること。
\item リアルタイムで発生するマイニング操作を通じて新しいビットコインを作成すること。
\end{enumerate}

\bigskip


\subsection*{Question 233}
Lightning Networkのマイクロトランザクションにおいて、以下のうちキーとなる機能ではないものはどれですか？

\begin{enumerate}[label=\Alph*., leftmargin=*]
\item 即時支払い。
\item 長い支払い時間。
\item スケーラビリティ。
\item 低い手数料。
\end{enumerate}

\bigskip


\subsection*{Question 234}
Lightning Networkがコンテンツクリエーターにもたらす可能性のある利点の一つは何ですか？

\begin{enumerate}[label=\Alph*., leftmargin=*]
\item ユーザーの注目を保持するために、より迅速にコンテンツを制作できるようになります。
\item 彼らはコンテンツに対してより多くの料金を請求できるようになります。
\item より広い視聴者に到達することができるようになります。
\item ユーザーの注目を保持するために、良質なコンテンツを制作することへのインセンティブが与えられます。
\end{enumerate}

\bigskip


\subsection*{Question 235}
DCAプラットフォームとは何ですか？

\begin{enumerate}[label=\Alph*., leftmargin=*]
\item 投資機能なしでリアルタイムに株価を追跡するソフトウェアツールです。
\item ビットコインの取引のみを許可する暗号通貨交換の一種です。
\item 投資家にマージントレーディングのためのローンを提供する金融サービスです。
\item 定期的な間隔で一定額のお金を資産に投資することを可能にするサービスです。
\end{enumerate}

\bigskip


\subsection*{Question 236}
Lightning Networkを利用して、物品のレンタルにどのように使用できる可能性がありますか？

\begin{enumerate}[label=\Alph*., leftmargin=*]
\item ユーザーは、物品を使用した時間に応じて、分単位や秒単位で料金を請求されることがあります。
\item ユーザーは、物品の状態に基づいて料金を請求されることがあります。
\item ユーザーは、物品のレンタルに高額な料金を請求されることがあります。
\item ユーザーは、物品の価値に基づいて料金を請求されることがあります。
\end{enumerate}

\bigskip


\subsection*{Question 237}
Lightning Networkによって提供される潜在的なビジネスチャンスは何ですか？

\begin{enumerate}[label=\Alph*., leftmargin=*]
\item 消費者がパーソナライズされたブロックチェーンを使用する新しい経済モデル。
\item 消費者がどんなコンテンツにも支払わない新しい経済モデル。
\item 消費者が消費した内容に基づいて支払う新しい経済モデル。
\item 消費者が異なる暗号通貨を使用して支払う新しい経済モデル。
\end{enumerate}

\bigskip


\subsection*{Question 238}
Lightning Networkを自分で試す一つの方法は何ですか？

\begin{enumerate}[label=\Alph*., leftmargin=*]
\item ポッドキャストアプリケーション「Fountain」を試すことです。
\item 新しい暗号通貨を作成すること。
\item 個人のコンピューターを使用してビットコインをマイニングすること。
\item 銀行からビットコインを購入すること。
\end{enumerate}

\bigskip


\subsection*{Question 239}
AI技術の進歩の一つの潜在的な影響は何ですか？

\begin{enumerate}[label=\Alph*., leftmargin=*]
\item AI技術は、今日の仕事の80％以上を消滅させるでしょう。
\item AI技術は世界経済の崩壊を引き起こすでしょう。
\item AI技術は人間の知能を減少させるでしょう。
\item AI技術はすべてのコンピュータを破壊するでしょう。
\end{enumerate}

\bigskip


\subsection*{Question 240}
ビットコインが技術と社会の未来にどのような関連があるのか？

\begin{enumerate}[label=\Alph*., leftmargin=*]
\item ビットコインは、中央サーバーを通じて情報を交換することを可能にするAI技術の一形態です。
\item ビットコインは、世界の人口を制御するツールであり、貧困の撲滅を可能にします。
\item ビットコインは、信頼できる第三者なしに価値を交換することを可能にする技術革命です。
\item ビットコインは、政府が市民の金融取引を追跡する方法です。
\end{enumerate}

\bigskip


\subsection*{Question 241}
金融の未来に関して提起される主な疑問の一つは何ですか？

\begin{enumerate}[label=\Alph*., leftmargin=*]
\item お金を完全になくす方法は？
\item 私たちが使用するお金を誰が保持し、認証し、追跡すべきか？
\item お金を時代遅れにする方法は？
\item すべての人が同じ額のお金を持つようにするにはどうすればいいですか？
\end{enumerate}

\bigskip


\subsection*{Question 242}
Bitcoinの主な特徴の一つは何ですか？

\begin{enumerate}[label=\Alph*., leftmargin=*]
\item 選ばれた人々によってのみ使用できます。
\item どのような権限もなしに世界中どこでも交換することができます。
\item 特定の国でのみ使用できます。
\item 使用するためには中央権威からの認可が必要です。
\end{enumerate}

\bigskip


\subsection*{Question 243}
ビットコインのような新しい技術を受け入れることの潜在的な利点の一つは何ですか？

\begin{enumerate}[label=\Alph*., leftmargin=*]
\item 世界中で大規模な規模の経済を生み出す可能性があります。
\item 世界的な経済崩壊を引き起こす可能性があります。
\item すべての仕事がなくなる可能性があります。
\item 人間の知能が低下する可能性があります。
\end{enumerate}

\bigskip


\subsection*{Question 244}
銀行システムに関して主な懸念事項は何でしょうか？

\begin{enumerate}[label=\Alph*., leftmargin=*]
\item 銀行のルールが透明ではなく、すべての人にとって理解しやすいものではない。
\item 銀行システムは理解しやすすぎる。
\item 銀行システムは分散化されすぎている。
\item 銀行システムは透明すぎる。
\end{enumerate}

\bigskip


\subsection*{Question 245}
検閲に関連する主な問題の一つは何ですか？

\begin{enumerate}[label=\Alph*., leftmargin=*]
\item 表現の自由や集会の権利を破壊する可能性があります。
\item イノベーションの減少につながる可能性があります。
\item 表現の自由の強制につながる可能性があります。
\item 集会の権利への支持が増える可能性があります。
\end{enumerate}

\bigskip


\subsection*{Question 246}
銀行口座を持たない人々にとってのビットコインの主な利点の一つは何ですか？

\begin{enumerate}[label=\Alph*., leftmargin=*]
\item ビットコインは、最も裕福なユーザーに報酬を与えることで、取引において不平等を生み出します。
\item ビットコインは、高い社会的地位を持つ人々にのみ利益をもたらします。
\item ビットコインは、社会的地位や政治的立場に関係なく、取引において平等を実現します。
\item ビットコインは、銀行口座と左派の政治的立場を持つ人々にのみ利益をもたらします。
\end{enumerate}

\bigskip


\subsection*{Question 247}
Bitcoinプロトコルの主な特徴の一つは何ですか？

\begin{enumerate}[label=\Alph*., leftmargin=*]
\item 特定のユーザーに偏見を持っています。
\item 選ばれた一部のユーザーのみに利益をもたらします。
\item それはすべてのユーザーに対して中立です。
\item 特定のユーザーを他のユーザーよりも優遇します。
\end{enumerate}

\bigskip


\subsection*{Question 248}
ビットコインを入手する実用的な方法は何ですか？

\begin{enumerate}[label=\Alph*., leftmargin=*]
\item 規制されたプラットフォームまたは規制されていないプラットフォームを通じて購入する。
\item 物理的な世界で見つける。
\item ラップトップを使用してマイニングする。
\item 特別なソフトウェアを使用して作成する。
\end{enumerate}

\bigskip


\subsection*{Question 249}
ビットコインが作られた主な理由の一つは何ですか？

\begin{enumerate}[label=\Alph*., leftmargin=*]
\item お金の廃止を提案すること。
\item 世界の人口の制御を提案すること。
\item 金融システムの変更を提案すること。
\item 新しい形のAI技術の創造を提案すること。
\end{enumerate}

\bigskip


\subsection*{Question 250}
ビットコインが銀行システムに与える主な影響の一つは何ですか？

\begin{enumerate}[label=\Alph*., leftmargin=*]
\item 銀行システムの排除を可能にします。
\item 銀行システムの腐敗を可能にします。
\item 銀行システムからの自立を可能にします。
\item 銀行システムの制御を可能にします。
\end{enumerate}

\bigskip


\newpage
\section*{Answer Key}

\begin{enumerate}
\item \textbf{Question 1:} C
\\ \textit{ブロックサイズ戦争は2017年に起こりました。この対立は、トランザクション容量を増やすためにブロックサイズを増加させることでBitcoinを変更したいと考えるアクターと、サイズを1MBに保つことでユーザーの独立性と権力を保持しようとしたアクターとの間で発生しました。}

\item \textbf{Question 2:} C
\\ \textit{このビットコインコースの目的は、ビットコインの一般的な金融特性について議論し、ビットコインの購入や売却、デジタルウォレットでの安全な保管、取引での使用方法などを含みます。また、ビットコインの将来について探求することも目的としています。}

\item \textbf{Question 3:} C
\\ \textit{Bitcoinは、暗号技術やブロックチェーンなどの技術を使用するコンピュータプロトコルであり、また、取引に使用される通貨単位でもあります。}

\item \textbf{Question 4:} A
\\ \textit{このビットコイン講座は、ビットコインの金融面に焦点を当てています。これには、ビットコインの購入や販売方法、デジタルウォレットでの安全な保管方法、取引での使用方法が含まれます。
レビュー済み: true}

\item \textbf{Question 5:} B
\\ \textit{ビットコインは中立的な通貨であり、どのような団体や機関にも制御されていないため、すべてのユーザーが同じ権利を持っています。これは、その分散化の重要な側面です。}

\item \textbf{Question 6:} A
\\ \textit{サイファーパンクスは、暗号技術が個人の権利を政府や大企業の干渉から守るためのツールとして機能すると信じた個人のグループです。}

\item \textbf{Question 7:} B
\\ \textit{Cypherpunkマニフェストは1993年にEric Hughesによって書かれ、プライバシーは基本的な権利であると主張しています。}

\item \textbf{Question 8:} A
\\ \textit{David Chaumは1980年代に、彼のプロジェクトDigiCashを通じて匿名電子マネーの概念を紹介しました。}

\item \textbf{Question 9:} C
\\ \textit{サイバースペースの独立宣言は1996年にジョン・ペリー・バーロウによって書かれました。この文書は、サイバースペースは物理的な領域とは異なる独自の領域であり、産業政府はそこで自らの権威を強制できないと主張しています。}

\item \textbf{Question 10:} A
\\ \textit{Wei Daiのb-moneyは、中央の実体がない匿名のデジタル通貨のアイデアであり、Bitcoinの創造に大きな影響を与えました。実際、Bitcoinのホワイトペーパーで引用されています。}

\item \textbf{Question 11:} B
\\ \textit{技術を超えて、ビットコインは伝統的な金融システムに対する革命を象徴しています：このプロトコルは、中央のエンティティなしに新しい通貨システムを強制します。}

\item \textbf{Question 12:} A
\\ \textit{最初の通貨形態は、穀物、家畜、その他の商品など、必需品に関連するものが多かったです。
しかし、これらの商品には、腐敗性という大きな欠点があり、長期的な貯蓄手段として使用することが困難でした。}

\item \textbf{Question 13:} A
\\ \textit{アリストテレスによると、お金は三つの機能を果たします：それは価値の保存、交換の媒体、そして計算の単位です。現在、交換の媒体機能が他の二つの前に、主に認識されています。}

\item \textbf{Question 14:} B
\\ \textit{効果的な通貨は、代替可能（価値の損失なく交換可能）、分割可能（異なるボリュームの取引を容易にするために小さな部分に分けられる）、そして流動性がある（容易に商品やサービスに変換可能）でなければなりません。}

\item \textbf{Question 15:} B
\\ \textit{金をお金の形態として使用する際の主な欠点の一つは、必要に応じて小さな単位に分割することが難しいことです。これは、使用する上で実用的ではないことを意味します。}

\item \textbf{Question 16:} D
\\ \textit{金が標準的な交換媒体となったのは、その希少性、耐久性、および分割可能性によるものです：その希少性は生産による価値の減少を防ぎ、その耐久性は時間が経っても腐食しないことを保証し、その分割可能性は異なるサイズの取引での使用を可能にしました。}

\item \textbf{Question 17:} A
\\ \textit{お金をコミュニケーションツールとして使用することで、私たちは自分の時間とエネルギーを将来にわたって価値のある資産に変換し、価値の減少のリスクなしに再利用できるようになります。また、お金は普遍的な共通言語としてのコミュニケーションを可能にし、見知らぬ二人が交換、取引を行い、物事の価値について合意することを可能にします。}

\item \textbf{Question 18:} D
\\ \textit{ドルやユーロなどのフィアット通貨の主な欠点は、その価値が通貨インフレーションによって常に侵食されることです。これは、それらを制御するエンティティ（王、中央銀行、皇帝、独裁者）による価値の減少が原因です。}

\item \textbf{Question 19:} D
\\ \textit{Bitcoinが他の通貨形態に対して持つ主な利点は、厳格に限定された供給量と信頼される第三者の不在により、優れた価値の保存手段を提供することです。これは長期的な貯蓄に適した選択肢となります。}

\item \textbf{Question 20:} A
\\ \textit{金は、分割や容易な交換を可能にする中央の実体が必要であるため、グローバル化されたデジタルの世界で追いつくことができません。これは、取引がますますデジタル化され、分散化される世界での使用が実用的でなくなることを意味します。}

\item \textbf{Question 21:} A
\\ \textit{厳格に限定された供給などの特性にもかかわらず、ビットコインは今日、商業で広く受け入れられていません。これは、個人による採用が不足しているためであり、その結果、商人がそれを受け入れることを奨励されていません。}

\item \textbf{Question 22:} A
\\ \textit{通貨は時間と共に進化する技術であり、最初は生の石で表され、その後、金属製のコイン、紙幣、銀行カードへと変化し、最終的にはビットコインのブロックチェーンとライトニングネットワークを使用したピアツーピアネットワークへと進化しました。}

\item \textbf{Question 23:} C
\\ \textit{金は密度が高く重いため、大量に持ち運ぶのは物理的に困難です。この重さは、より軽い素材と比較して輸送に便利ではないことを意味します。}

\item \textbf{Question 24:} D
\\ \textit{ドルやユーロなどのフィアット通貨は、現在ほとんどがデジタル化されているため、非常に流動性が高く、簡単に運ぶことができます。しかし、その価値は通貨インフレによって常に侵食されています。}

\item \textbf{Question 25:} A
\\ \textit{ビットコインは、分割可能性、検証可能性、許可不要の特性、持ち運びやすさなどの特性により、優れた価値の保存手段を提供します。さらに、中立的なインターネット通貨として、国境を知らない良好な交換手段を表しています。}

\item \textbf{Question 26:} B
\\ \textit{健全な通貨は、市場と自発的な合意から生まれなければならない自然な社会現象です。人間や人間の集団が通貨を作り出すことはできません。}

\item \textbf{Question 27:} C
\\ \textit{ビットコインは、中央の組織がないピアツーピアネットワークから生まれるため、新しい形のお金と見なすことができます。}

\item \textbf{Question 28:} D
\\ \textit{金は簡単に分割したり、長距離を運ぶことは難しいですが、非常に耐久性のある素材であるため、価値の保存手段として適しています。過去には金貨が通貨として使用されましたが、今日のデジタル時代では同じ目的を果たすことはありません。}

\item \textbf{Question 29:} A
\\ \textit{フィアットマネーは政府の布告によって法定通貨として確立され、その価値はどのような物理的な商品からも派生していません。}

\item \textbf{Question 30:} C
\\ \textit{Fiduciary通貨は、金や銀のような物理的な商品によって裏打ちされていません。その価値は、それらを発行し規制する機関への人々の信頼と信用から派生しています。}

\item \textbf{Question 31:} C
\\ \textit{中央銀行は信託通貨を発行する機関です。彼らは金融政策を決定し、どれだけのお金を流通させるか、または印刷するかを決めます。}

\item \textbf{Question 32:} A
\\ \textit{Bitcoinの総供給量は2100万単位に制限されています。
これはBitcoinを中央銀行が無制限に発行できる信託通貨と
区別する重要な特徴です。
レビュー済み: true}

\item \textbf{Question 33:} C
\\ \textit{Bitcoinは、国家とお金を分離することを目的として作られました。これは、国家がお金の供給を完全にコントロールし、自らの利益のためにそれを操作できる現行の金融システムの体系的な課題に対する反応です。}

\item \textbf{Question 34:} A
\\ \textit{中央機関は、初期のものと比較して、毎年数パーセント通貨量を増加させることにより、通貨の価値をしばしば下げます。この静かな価値下げは、しばしば人々の利益になると正当化されます。}

\item \textbf{Question 35:} C
\\ \textit{通貨印刷、または中央銀行による新しいお金の創造は、インフレを引き起こします。これは、通貨供給の増加がお金の購買力を減少させ、価格の上昇と貯蓄の価値の減少を引き起こすためです。}

\item \textbf{Question 36:} D
\\ \textit{中央銀行デジタル通貨（CBDC）は、個人の金融の自由を妨げ、権威主義的な乱用を容易にする可能性があります。これは、より中央計画的な経済を提供し、中央銀行が通貨供給をさらにコントロールすることになるためです。}

\item \textbf{Question 37:} B
\\ \textit{歴史的に、指導者たちは金を中央集権化し、金と同等の価値を持つ新しい通貨を導入しました。この新しい通貨は、しばしば人々の利益のためであるという名目の下で、徐々に価値が下げられています。}

\item \textbf{Question 38:} A
\\ \textit{信託通貨に対する突然の価値下落または信頼の喪失は、ハイパーインフレを引き起こす可能性があります。これは、商品やサービスの価格が極端に高い割合で増加し、しばしば経済の不安定につながる状況です。}

\item \textbf{Question 39:} B
\\ \textit{ビットコインは分散型であり、暗号技術に基づいているため、偽造することができず、21百万単位に限定されています。これは、中央銀行によって変更され得る通貨政策に従う法定通貨と異なる主要な特徴です。}

\item \textbf{Question 40:} B
\\ \textit{政治的アクターが現在の金融システムから恩恵を受けているため、彼らは根本的な変更を行う動機付けがされていません。
これにより、システムは自身の崩壊を避けるために規制され、制限されているにもかかわらず、可能な崩壊に至るまでそのまま続けられます。}

\item \textbf{Question 41:} D
\\ \textit{ハイパーインフレーションは、通貨発行によるインフレの急激な増加を特徴とする通貨現象であり、通貨への完全な信頼喪失につながります。}

\item \textbf{Question 42:} A
\\ \textit{ハイパーインフレのため、個人は迅速に通貨への信頼を完全に失い、これが短期間にわたってその価値の大幅な減少につながります。}

\item \textbf{Question 43:} D
\\ \textit{ハイパーインフレの最初の段階は、通貨への信頼の喪失です。}

\item \textbf{Question 44:} B
\\ \textit{ハイパーインフレの最終段階は、古いものを置き換えるために新しい通貨が導入されるときに発生します。
レビュー済み: true}

\item \textbf{Question 45:} D
\\ \textit{2%のインフレでは、毎年購買力が2%失われ、5年間で合計10%失われます。}

\item \textbf{Question 46:} D
\\ \textit{20%のインフレでは、3年間で購買力のほぼ半分を失います。}

\item \textbf{Question 47:} A
\\ \textit{1日につきインフレ率が100%というのは、実際に起こった現実的なシナリオです。それは2007年のジンバブエで起こり、お金の価値が約24時間ごとに倍になりました。}

\item \textbf{Question 48:} C
\\ \textit{ハイパーインフレの第二段階は、信頼喪失フェーズの直後に起こる通貨崩壊 & 物価上昇です。}

\item \textbf{Question 49:} D
\\ \textit{ドイツで最も高いインフレ率を記録したのは1923年10月で、インフレ率は29,500%に達しました。}

\item \textbf{Question 50:} B
\\ \textit{ハンガリーで最もインフレ率が高かった月は1946年7月で、記録的な価格インフレ率は41,900,000,000,000,000%でした。}

\item \textbf{Question 51:} D
\\ \textit{2008年11月、ジンバブエの通貨の月間インフレ率は79,600,000,000%に達し、ジンバブエドルの価値を0にまで激減させました。}

\item \textbf{Question 52:} B
\\ \textit{2009年4月、財務大臣はジンバブエドルの使用停止を発表し、貿易のために異なる外貨の使用を認可しました。}

\item \textbf{Question 53:} D
\\ \textit{ビットコインは、効率的で健全な通貨であるための必要な特徴をすべて備えています：分割可能、即時に輸送可能、検閲不可能、検証コストが無視でき、そして21百万単位で数世紀にわたって設定された金融政策を持っています。}

\item \textbf{Question 54:} D
\\ \textit{Bitcoinネットワーク上で取引を検証するプロセスはマイニングと呼ばれています。マイナーは暗号学的なタスクを実行し、新しいビットコインの発行によって報酬を得ます。}

\item \textbf{Question 55:} A
\\ \textit{マイニング報酬が半分になるイベントはハービングと呼ばれます。このイベントにより、通貨発行のカーブは階段のような形状をしています。}

\item \textbf{Question 56:} C
\\ \textit{ブロックチェーンに新しいブロックが10分ごとに追加されるメカニズムは難易度調整です。
この調整は、約2016ブロックごと、または約2週間ごとに行われます。}

\item \textbf{Question 57:} C
\\ \textit{ゲーム理論は、人間の合理性に基づいた数学的概念です。ビットコインでは、プロトコルの変更はユーザーによって強制されるため、ビットコインプロトコルへの任意の変更を行うことは非常に複雑です。
レビュー済み: true}

\item \textbf{Question 58:} A
\\ \textit{Bitcoinノード上で流通しているビットコインの量を確認するコマンドはbitcoin-cli gettxoutsetinfoです。このコマンドは透明性と検証可能性を提供し、Bitcoinシステムへの信頼を強化します。
レビュー済み: true}

\item \textbf{Question 59:} C
\\ \textit{ビットコインはオーストリア経済学の原則に沿っています。
その上限数量と中央機関の欠如は、伝統的な通貨で見られる
インフレーションの固有のリスクからそれを保護します。}

\item \textbf{Question 60:} B
\\ \textit{ハービングメカニズムにより、ビットコインの生成が2140年に停止されると数学的に予測できます。この時点でビットコインの総数がその上限である2100万に達するためです。
レビュー済み: true}

\item \textbf{Question 61:} C
\\ \textit{難易度調整は、グローバルなハッシュレートに関係なく、平均して10分ごとに新しいブロックがブロックチェーンに追加されるようにするため、約2週間または2016ブロックごとに行われるメカニズムです。
レビュー済み: true}

\item \textbf{Question 62:} B
\\ \textit{マイニングの報酬は、210,000ブロックごとに半減するようにプログラムされています。これはおよそ4年ごとのことで、「ハービング」と呼ばれるイベントです。}

\item \textbf{Question 63:} C
\\ \textit{最初の半減期はブロック高さ210,000で発生し、その後の流通しているビットコインの推定数は10,500,000 BTCでした。これは、流通する予定のビットコインの総数の半分です。
レビュー済み: true}

\item \textbf{Question 64:} B
\\ \textit{第4回半減期は、ブロック高さ840,000で発生し、その後のマイニングでのBTC報酬は3.125 BTCになります。}

\item \textbf{Question 65:} A
\\ \textit{7回目の半減期は2036年に起こる予定で、その時点でのビットコインの流通量は20,835,937.5になると推定されており、これは総供給量の約99.22%に相当します。
レビュー済み: true}

\item \textbf{Question 66:} A
\\ \textit{ビットコインウォレットはビットコインネットワークとのやり取りに使用されます。その3つの主な機能は、ビットコインを受け取ることを可能にする、ビットコインを送ることを可能にする、そしてハッキングや盗難の試みからそれらを守ることです。}

\item \textbf{Question 67:} A
\\ \textit{Bitcoinウォレットを初期化する際に、ニーモニックシードと呼ばれる秘密の単語リストが生成されます。 この回復リストは非常に重要で、Bitcoinの所有権を表し、それゆえにそれらを送信する権利を意味します。}

\item \textbf{Question 68:} A
\\ \textit{プライベートキーはあなたのウォレットに保存されていますが、ビットコイン自体は実際にはビットコインブロックチェーンに保存されています。ビットコインブロックチェーンとは、ビットコインのピアツーピアネットワーク内の公開分散台帳のことです。}

\item \textbf{Question 69:} A
\\ \textit{2017年以降、プライベートキーは12または24の単語のシンプルなリスト、いわゆるニーモニックフレーズにコード化することができます。このフレーズはビットコインウォレットのバックアップであり、任意のビットコインウォレットソフトウェア/アプリを使用してビットコインにアクセスすることを可能にします。}

\item \textbf{Question 70:} D
\\ \textit{公開鍵は秘密鍵から派生され、お金を受け取るためにBitcoinアドレスを生成するために使用されます。}

\item \textbf{Question 71:} A
\\ \textit{公開鍵を共有することは、セキュリティではなくプライバシーにリスクをもたらします。これは、公開鍵がBitcoinアドレスの生成に使用され、その結果お金を受け取るために使われるが、Bitcoinの所有権を表すものではないためです。}

\item \textbf{Question 72:} C
\\ \textit{ウォレットを保持しているデバイスを失ったからといって、必ずしもBitcoinを失うわけではありません。ウォレットを再作成し、Bitcoinを使用することを可能にするのは、ニーモニックフレーズです。}

\item \textbf{Question 73:} C
\\ \textit{ビットコインウォレットのデフォルトのセキュリティに満足できない場合は、ビットコインウォレットにパスフレーズを追加することで、いつでもセキュリティを強化することができます。}

\item \textbf{Question 74:} A
\\ \textit{ウォレットを作成するために使用される暗号技術のおかげで、12または24の単語のリストを誰かが偶然推測する確率は極めて低いです。これは、1から$2^256$の間の「正しい」数値を見つけることとほぼ同等であり、宇宙で定義された原子を見つけることに相当します。}

\item \textbf{Question 75:} D
\\ \textit{ビットコインウォレットを選択する際の中心的な質問は、資金の所有者はあなた自身ですか、それともお金の管理を第三者に委ねていますか？ということです。これは、プライベートキーの保持者がビットコインの所有者であるためです。}

\item \textbf{Question 76:} D
\\ \textit{Bitcoinを使用するには、所有者が楕円曲線デジタル署名アルゴリズム(ECDSA)を使用して署名を生成し、Bitcoinトランザクションの伝播のためにこれを用います。}

\item \textbf{Question 77:} D
\\ \textit{ビットコインを送信者にあなたの資金を盗ませることなく受け取る方法を説明するために使用されている比喩は郵便受けです。人々がお金を預けますが、それを開けることができるのはあなただけです。}

\item \textbf{Question 78:} B
\\ \textit{ビットコインを所有している場合、主な懸念事項は資金のセキュリティです。これは、ビットコインを含む他の形態の通貨と同様に、適切に保護されていない場合に盗まれたり失われたりする可能性があるためです。}

\item \textbf{Question 79:} A
\\ \textit{ホットウォレットとは、インターネットアクセスが可能なデバイス、例えば携帯電話やコンピューターに保存されたビットコインウォレットを指す用語です。この方法では、取引を簡単に行うことができますが、ウォレットがオンラインの脅威にさらされる可能性も意味しています。}

\item \textbf{Question 80:} D
\\ \textit{コールドウォレットとは、インターネットに接続されていないデバイスに保存されるBitcoinウォレットを指す用語です。これにより、オンラインの脅威に対する脆弱性が低減されますが、日常の取引においてウォレットがアクセスしにくくなるという意味もあります。}

\item \textbf{Question 81:} C
\\ \textit{マルチシグウォレットは、トランザクションを実行するために複数の署名が必要なビットコインウォレットの一種です。これは、単一の人物が他の人々の承認なしにトランザクションを実行できないことを意味し、追加のセキュリティ層を提供します。}

\item \textbf{Question 82:} C
\\ \textit{カストディアルサービスを使用してビットコインを保管する場合、あなたはビットコインの唯一の保持者ではありません。これは、信頼できる第三者がいつでもあなたの資金へのアクセスを制限することができるということを意味し、伝統的な銀行システムと同じレベルの財政的主権をあなたに与えます。}

\item \textbf{Question 83:} A
\\ \textit{ビットコインのセキュリティとアクセシビリティを複雑にしすぎると、例えば、ウォレットのバックアップを誤って取り扱った場合に自分自身に害を及ぼす可能性があります。これは、携帯電話やコンピュータの紛失の場合に資金へのアクセスを再度得るためには、シンプルなバックアップ手順が重要であるためです。}

\item \textbf{Question 84:} A
\\ \textit{モバイルウォレットは、取引を簡単に行えるため、日常の支出に推奨されます。しかし、オンラインの脅威に対して潜在的に脆弱であるため、大量の金額や長期間の貯蓄には推奨されません。
レビュー済み: true}

\item \textbf{Question 85:} C
\\ \textit{オフライン、またはコールドの物理ウォレットは、オンラインの脅威に対してホットウォレットよりも脆弱性が低いため、大量の保管に推奨されます。しかし、日常の取引にはアクセスしにくいため、日常の支出には推奨されません。}

\item \textbf{Question 86:} D
\\ \textit{マルチシグウォレットは、トランザクションを実行するために複数の署名が必要とされるため、大量の保管および企業用途に推奨されます。このプロセスは、単一の人物が他の人々の承認なしにトランザクションを実施できないため、追加のセキュリティ層を加えます。}

\item \textbf{Question 87:} B
\\ \textit{パスフレーズ付きのウォレットを使用する場合、携帯電話やコンピュータの紛失時に資金へのアクセスを再度得るためには、12または24の単語のリストとパスフレーズの両方をバックアップする必要があります。一方または他方がなければ、資金にアクセスすることはできません。}

\item \textbf{Question 88:} D
\\ \textit{マルチシグウォレットを使用する場合、各プライベートキーは異なる場所に保管されるべきです。これは、マルチシグがトランザクションを実行するために複数の署名を必要とするため、必要な数のウォレットが必要になるからです。したがって、各部分を異なる場所に保管することで、セキュリティの追加の層が加わります。}

\item \textbf{Question 89:} B
\\ \textit{用語「hodl」は、2013年にビットコインフォーラムのオンライン投稿で誤って綴られたことから起源を持ちます。時間が経つにつれて、「hodl」はより広い投資戦略を代表するように進化しました。市場の変動にかかわらず、長期間にわたって暗号通貨を保持し続けることを投資家に奨励します。}

\item \textbf{Question 90:} D
\\ \textit{Bitcoinウォレットでは、プライベートキーは通常、12または24の単語のリスト、別名「シードフレーズ」または「ニーモニックフレーズ」として表されます。}

\item \textbf{Question 91:} A
\\ \textit{私的な鍵が第三者に明らかにされた場合、関連する資金はもはや安全ではありません。なぜなら、私的な鍵を通じて、彼らはあなたの資金にアクセスできるからです。}

\item \textbf{Question 92:} C
\\ \textit{ニーモニックフレーズを紙に保存する代わりに推奨される方法の一つは、金属板に刻むことです。この解決策はより耐久性があり、長持ちします。}

\item \textbf{Question 93:} C
\\ \textit{ニーモニックフレーズの単語を書き留めた後、二つ目のコピーを作成し、最初のコピーとは別の場所に保管することをお勧めします。これにより、最初のコピーが紛失したり事故が起きたりした場合のバックアップが提供されます。
レビュー済み: true}

\item \textbf{Question 94:} D
\\ \textit{ウォレットに12語または24語のリスト、つまりニーモニックフレーズがない場合、標準的なバックアップ手順に従っていないため、ウォレットが安全でないことに注意すべきです。}

\item \textbf{Question 95:} C
\\ \textit{Bitcoinウォレットのプライベートキーをバックアップする標準的な方法は、任意のウォレットソフトウェアにニーモニックフレーズを入力することです。これにより、ウォレットを復元できます。
レビュー済み: false}

\item \textbf{Question 96:} C
\\ \textit{鋼のような耐久性のある素材にニーモニックフレーズを刻むことで、水損害や火災にも耐えることができるキーの物理的なバックアップを作成し、長期的にビットコインを保護します。
レビュー済み: true}

\item \textbf{Question 97:} C
\\ \textit{相続計画により、死後にビットコインが適切に管理されることが保証されます。これには、資産の詳細、アクセス方法、連絡先情報を記載した手書きの手紙や、連絡を取るべき信頼できる個人の連絡先情報が含まれることがあります。}

\item \textbf{Question 98:} C
\\ \textit{Bitcoinの文脈におけるニーモニックフレーズとは、12または24の単語のリストであり、これにより任意のBitcoinウォレットアプリケーションであなたのウォレットを復元することができます。このリストにアクセスできる人は誰でもあなたのビットコインにアクセスすることができます。}

\item \textbf{Question 99:} C
\\ \textit{ビットコインの文脈でプライバシーが重要なのは、資金の追跡リスクを最小限に抑えることで、監視を容易にすることを防ぐためです。}

\item \textbf{Question 100:} A
\\ \textit{ビットコインを取引プラットフォームに残しておくことは推奨されません。なぜなら、ハッカー攻撃の対象となりやすいからです。}

\item \textbf{Question 101:} A
\\ \textit{Bitcoin相続の文脈において、ビットコインの相続について公証人と話し合うことは、税務コンプライアンスを確保するために重要です。}

\item \textbf{Question 102:} B
\\ \textit{ビットコインウォレットは、ビットコインを取引するために必要なプライベートキーを保存するために使用されるソフトウェアの一種です。}

\item \textbf{Question 103:} D
\\ \textit{ビットコインでは、金融の主権は個人の責任と密接に関連しています。なぜなら、推奨されるバックアップ手順に従って自分の資金を自己管理する必要があるからです。}

\item \textbf{Question 104:} C
\\ \textit{Blockmitは、鋼のような耐久性のある素材にニーモニックフレーズを刻むための低コストの解決策であり、
長期的にビットコインを安全に保管することができます。}

\item \textbf{Question 105:} B
\\ \textit{KYCまたは顧客を知ることは、企業がクライアントの身元を確認するために採用する手続きを指します。
これは通常、規制への準拠を確保し、マネーロンダリングや詐欺などの違法行為を防ぐために行われます。}

\item \textbf{Question 106:} B
\\ \textit{KYC非対応のビットコインを購入するとは、多くの取引所やプラットフォームが要求する標準的な「顧客確認」(Know Your Customer) 検証手続きを経ずにビットコインを取得することを意味します。個人がこの方法を選択する主な理由の一つは、個人情報や財務活動を公開せずにプライバシーを維持することです。}

\item \textbf{Question 107:} C
\\ \textit{一人の人がセキュリティを優先する一方で、別の人はユーザーフレンドリーさや特定の機能を重視するかもしれません。したがって、単一のBitcoinウォレットが全ての人の要件を満たすことはなく、複数のウォレットオプションが存在する必要があります。}

\item \textbf{Question 108:} D
\\ \textit{ビットコインが初めて世界に紹介されたのは2008年10月31日で、サトシ・ナカモトがビットコインのホワイトペーパーを含むメールをサイファーパンクのメーリングリストに送信したときでした。}

\item \textbf{Question 109:} D
\\ \textit{Hal FinneyはSatoshi Nakamoto自身の後でBitcoinネットワークに参加した最初の人物でした。彼はその機会を記念してRunning Bitcoinとツイートしました。
レビュー済み: true}

\item \textbf{Question 110:} A
\\ \textit{最初に記録されたBitcoin取引は、Satoshi NakamotoとHal Finneyの間の10 BTCの取引でした。}

\item \textbf{Question 111:} C
\\ \textit{サトシ・ナカモトは2010年に姿を消し、彼の創造物をコミュニティの手に委ねました。}

\item \textbf{Question 112:} B
\\ \textit{ビットコインブロックチェーンの最初のブロック、別名
ジェネシスブロックには、03/jan/2009
銀行への二度目の救済措置の瀬戸際にいる財務大臣というメッセージが含まれています。これは
ジェネシスブロックのタイムスタンプとしてその日のタイム誌のタイトルです。}

\item \textbf{Question 113:} B
\\ \textit{BitcoinTalkフォーラムは、Bitcoinユーザーのための議論の場です。これは、2009年11月22日にサトシ・ナカモトによって作成されました。}

\item \textbf{Question 114:} B
\\ \textit{ビットコインが初めて物理的な商品の購入に使われたのは、2010年5月22日にLaszlo Hanyeczが10,000 BTCで2枚のピザを購入したときで、これは現在ピザの日として知られています。}

\item \textbf{Question 115:} A
\\ \textit{ビットコインのオンラインでの可用性は、2009年以来99.988%です。}

\item \textbf{Question 116:} D
\\ \textit{2009年以降、ビットコインのオンラインでの可用性は99.988%です。}

\item \textbf{Question 117:} C
\\ \textit{最初のBitcoin取引は、Satoshi NakamotoとHal Finneyの間で行われた10BTCの取引で、ブロック170に記録されています。}

\item \textbf{Question 118:} C
\\ \textit{最初にリリースされたBitcoinソフトウェアはBitcoin-0.1.0でした。これは2009年1月8日にサトシ・ナカモトによって発表されました。 レビュー済み: true}

\item \textbf{Question 119:} B
\\ \textit{サトシ・ナカモトがメールで送った最後の知られているプライベートメッセージの日付は2011年4月23日です。このメッセージは開発者のマイク・ハーンに送られました。}

\item \textbf{Question 120:} D
\\ \textit{ビットコイン取引は、ビットコインの所有権をデジタルで移転することです。これは、楕円曲線暗号技術のおかげで導出されたビットコインアドレスを使用して行われます。このアドレスは、ビットコインを受け取る人のウォレットに固有のものです。}

\item \textbf{Question 121:} D
\\ \textit{ビットコイン取引における取引手数料は、マイナーに次のブロックに取引を含めるためのインセンティブとして機能します。これらの手数料は、ブロックに取引を含めるための自由市場を生み出します。
レビュー済み: true}

\item \textbf{Question 122:} B
\\ \textit{相続計画を作成することは、あなたが亡くなった後にビットコインが適切に管理されることを確実にするための重要なステップです。
この計画には、資産のリストアップした手書きの手紙、それらのアクセス方法、および連絡されるべき信頼できる個人の連絡先情報を含めることができます。}

\item \textbf{Question 123:} D
\\ \textit{ビットコインのブロックチェーンは、すべてのビットコイン取引の公開されていて変更不可能な台帳です。これはすべてのビットコインユーザーの共通の台帳です。}

\item \textbf{Question 124:} B
\\ \textit{ビットコイン取引における確認数は、その取引を含むブロックの上に採掘されたブロックの数を示しています。取引の確認数が多いほど、その取引はより不変性を持ちます。}

\item \textbf{Question 125:} B
\\ \textit{BitcoinネットワークにおけるBitcoinノードは、新しいブロックとその中のトランザクションの検証を扱います。このBitcoinノードのネットワークがあることで、真に分散化されたものとなっています。}

\item \textbf{Question 126:} A
\\ \textit{Bitcoinノードの地理的分布は、ブロックチェーンの全てのコピーを破壊することが実質的に不可能になります。ノードが広がっているため、物理的に全てを押収することは困難です。}

\item \textbf{Question 127:} B
\\ \textit{マイナーは「プルーフ・オブ・ワーク」と呼ばれるプロセスを通じて取引を検証し、ブロックに含めます。これには、暗号パズルを解く作業が含まれます。 レビュー済み: true}

\item \textbf{Question 128:} B
\\ \textit{Bitcoinネットワークのブロックマイニングの文脈において、有効な「ハッシュ」とは、ブロックに関連付けられた、256文字から成るユニークな指紋のことを指します。このハッシュの有効性は、Bitcoinネットワークの難易度に依存します。}

\item \textbf{Question 129:} A
\\ \textit{Bitcoinネットワークにおける難易度調整は、ブロックが約10分ごとに追加されるように保証します。これはBitcoinネットワークのプロトコルルールの一部です。}

\item \textbf{Question 130:} D
\\ \textit{Bitcoinネットワークは、プロトコルルールと経済的インセンティブを通じて、Bitcoinプロトコル内で悪意のある行為を行うコストを高くします。これにより、Bitcoinの所有権移転のどの段階でも信頼が不要になります。}

\item \textbf{Question 131:} B
\\ \textit{ビットコインネットワークのユーザーは、私的な鍵で取引にデジタル署名をすることによって、自分のお金の所有権を移転します。これにより、彼らが移転したいビットコインの所有者であることが確認されます。
レビュー済み: true}

\item \textbf{Question 132:} D
\\ \textit{Bitcoinネットワークのノードは、Bitcoinブロックチェーンのコピーを保持し、トランザクションを検証し、他のノードに情報を伝達し、Bitcoinプロトコルのルールを施行するなど、いくつかの重要な機能を果たします。}

\item \textbf{Question 133:} A
\\ \textit{2023年9月時点で、世界中には約45,000のノードが分散しています。 レビュー済み: true}

\item \textbf{Question 134:} B
\\ \textit{2023年時点で、Bitcoinブロックチェーンのサイズは約500GBです。}

\item \textbf{Question 135:} D
\\ \textit{2023年現在、Raspberry Pi 4でBitcoinノードを運用する際の電気消費による年間コストは10ユーロ弱です。
レビュー済み: true}

\item \textbf{Question 136:} A
\\ \textit{Bitcoinのコンテキストにおけるプルーニングされたノードは、メモリ内に最後のNブロックのみを保持するノードのことを指します。}

\item \textbf{Question 137:} C
\\ \textit{もしノードユーザーが現在の合意ルールを無効にする他のルールに従うことを決定した場合、そのノードはもはやBitcoinネットワークの一部ではなくなります。}

\item \textbf{Question 138:} C
\\ \textit{1ブロックが1MBで、10分ごとに生成されると考えると、Bitcoinノードを運用するためには、約5GB/月の帯域幅が必要です。 レビュー済み: true}

\item \textbf{Question 139:} A
\\ \textit{2017年のブロック戦争は、トランザクション容量を拡大するためにビットコインを修正し、ブロックサイズを増やすかどうかについての対立でした。}

\item \textbf{Question 140:} D
\\ \textit{ビットコインノードの手頃なコストとアクセシビリティは、ネットワークの分散化を促進するために重要です。}

\item \textbf{Question 141:} B
\\ \textit{もしブロックが100倍重くなった場合、Bitcoinノードを運用するには50TBのハードディスク、月に500GB以上の帯域幅、そして10分以内に数十万のトランザクションを検証できるハードウェアが必要になります。これにより、Bitcoinノードの運用が一般の人にとってアクセスしにくくなり、プロトコルの分散化とトランザクションおよび合意ルールの不変性が損なわれます。}

\item \textbf{Question 142:} A
\\ \textit{2017年の「ブロック戦争」では、ユーザーとノードはブロックサイズを増加させる提案を拒否し、SegWitと呼ばれるアップデートを活性化することで勝利しました。}

\item \textbf{Question 143:} C
\\ \textit{ライトニングネットワークは、ビットコインブロックチェーンに基づいた即時ビットコイン支払いネットワークです。}

\item \textbf{Question 144:} B
\\ \textit{マイナーは複雑な数学的問題を解くために計算能力を使用します。これにより、新しいトランザクションをブロックチェーンに追加することができます。このプロセスにより、ネットワークが保護され、その完全性が維持されます。}

\item \textbf{Question 145:} C
\\ \textit{Proof of Workは、取引を検証し新しいブロックを作成するためにBitcoinネットワークで使用される合意アルゴリズムです。マイナーが複雑な数学的問題を解決することを要求し、ネットワークのセキュリティと完全性を保証します。}

\item \textbf{Question 146:} A
\\ \textit{ハッシュレートは、マイナーがプルーフ・オブ・ワークを解決するために1秒間に行う試行の総数を指します。これは、ビットコインネットワークの計算能力を測る指標です。}

\item \textbf{Question 147:} B
\\ \textit{Bitcoinマイニングにおける難易度調整は、ネットワークの総計算能力に関係なく、ブロックをマイニングするのにかかる時間がおおよそ一定に保たれるように設計されています。これは、マイナーが解決しなければならない数学的な問題の難易度を、ネットワークに参加しているマイナーの数に基づいて調整します。}

\item \textbf{Question 148:} B
\\ \textit{ASICs（Application-Specific Integrated Circuits）は、
Bitcoinマイニングのために特別に設計された専門のハードウェアです。これらは、マイナーがProof of Workプロセスで解決しなければならない数学的問題であるSHA256ハッシュ関数を実行するために最適化されています。}

\item \textbf{Question 149:} C
\\ \textit{コインベーストランザクションは常にブロックの最初のトランザクションです。これには、検証者の作業を行ったマイナーへの報酬が含まれます。}

\item \textbf{Question 150:} C
\\ \textit{マイニングプールは、計算能力を組み合わせることでブロックのマイニングのチャンスを増やすマイナーのグループです。リソースをプールすることで、彼らはより頻繁にブロックをマイニングし、より小さいがより定期的な報酬を受け取ることができます。}

\item \textbf{Question 151:} B
\\ \textit{ビザンチン将軍問題は、分散コンピューティングにおける古典的な問題であり、一部のノードが故障しているか悪意がある可能性があるネットワークで合意に達することを含みます。ビットコインは、Proof of Workに基づいたブロックチェーンを使用することで、この問題を解決し、すべてのノードがトランザクション履歴の状態について合意することを保証します。}

\item \textbf{Question 152:} B
\\ \textit{Proof of Workは、ビットコインネットワークで取引を検証し、新しいブロックを作成するために使用される合意形成アルゴリズムです。マイナーが複雑な数学的問題を解く必要があり、これにより信頼できる第三者なしでネットワーク内の情報の真実性が保証されます。}

\item \textbf{Question 153:} C
\\ \textit{ビットコインマイニングにおけるエネルギーコストは、ネットワークのセキュリティにとって重要な側面です。マイニングにかかる高いエネルギーコストは、悪意のある行為者がブロックチェーンを操作することを極めて高価にします。一方で、取引の検証コストが低いことは、誰でもネットワークに参加できることを保証します。}

\item \textbf{Question 154:} C
\\ \textit{51%攻撃では、ネットワークのハッシュレートの半分以上を持つエージェントが、代替のブロックチェーンを作成することによってブロックチェーンを操作する可能性があります。
しかし、そのような攻撃は非常に高額なコストがかかり、利益を出すことが難しいため、起こりにくいシナリオです。}

\item \textbf{Question 155:} B
\\ \textit{ビットコインマイニングの主な環境上の懸念は、エネルギー消費です。ASICは主にアルミニウムで作られており、リサイクルまたは再利用が可能ですが、これらのマシンを動かすために必要なエネルギーは莫大です。}

\item \textbf{Question 156:} A
\\ \textit{Bitcoinプロトコルにおけるマイナーは、現在および過去のネットワークの両方を保護します。Bitcoinの価格を調整したり、その供給を制御したり、ユーザー間の取引を管理するわけではありません。}

\item \textbf{Question 157:} B
\\ \textit{Bitcoinプロトコルでは、平均して10分ごとにブロックが生成されます。この時間は、マイナーの数や計算能力に関わらず一定です。 レビュー済み: true}

\item \textbf{Question 158:} C
\\ \textit{ビットコインは検閲や銀行の制限から逃れることによって、金融の自由を提供します。これは特に、金融的抑圧や独裁政権下に住む個人にとって有益です。}

\item \textbf{Question 159:} A
\\ \textit{Bitcoinプロトコルは、2016ブロックごとに有効なハッシュを見つける難易度を調整することで、2つのブロック間の平均時間が一定に保たれるようにしています。このイベントは、マイナーの数や彼らの計算能力に関係なく発生します。}

\item \textbf{Question 160:} D
\\ \textit{鉱夫は、まだグリッドに接続されていない発電所のために、最後の手段としての買い手として行動することができます。これにより、発電所は電気ネットワークに接続される前にも資金調達を確保することができます。}

\item \textbf{Question 161:} C
\\ \textit{現在の金融システムは、過剰消費と借金を助長していると批判されており、これが深刻な環境問題を引き起こしています。これは、クレジットへの容易なアクセス、銀行による通貨発行、および部分準備銀行制度の使用によるものです。}

\item \textbf{Question 162:} B
\\ \textit{ビットコインはグリーンエネルギーの使用を促進し、それが生態系の移行を助ける可能性があります。例えば、汚染を防ぐためにメタンを燃やしている油井の炎は、ビットコインマイナーによって消火されることがあります。 レビュー済み: true}

\item \textbf{Question 163:} A
\\ \textit{ビットコインプロトコルは、次のブロックの有効なハッシュを見つけたマイナーに、合意ルールによって定義された報酬と、有効なブロックに含まれるすべてのトランザクションからの手数料を合わせた報酬を与えます。}

\item \textbf{Question 164:} B
\\ \textit{専用SHA256マシン、またはASICは、Bitcoinマイニングにおいてジュールあたりのハッシュレートを増加させます。これは、消費されたエネルギーあたりの秒間試行回数を意味します。}

\item \textbf{Question 165:} D
\\ \textit{ビットコインプロトコルが検閲耐性と回復力を保証するのは、プロトコルの各コンポーネントが地球上の地理的に分散された場所に配置されているためです。例えば、全大陸にわたって約40,000のビットコインノードが存在します。}

\item \textbf{Question 166:} A
\\ \textit{金融システムに関連するケインズ経済学とオーストリア経済学の考え方の主な違いは、ケインズ経済学が状況や資源の時間的および動的な側面を考慮に入れていないことです。言い換えれば、無制限の通貨は、私たちの惑星の限られた資源を効果的に反映することはできません。}

\item \textbf{Question 167:} C
\\ \textit{ビットコインの分散型の性質、匿名性、検閲への抵抗力は、テクノフィル、サイファーパンク、リバタリアン、金愛好家などのグループに魅力的です。}

\item \textbf{Question 168:} A
\\ \textit{Bitcoinをマイニングするには大量の計算能力が必要であり、それに伴って電力を消費します。 また、マイナーは効果的にマイニングを行うために、特殊なハードウェアへの投資が必要です。}

\item \textbf{Question 169:} A
\\ \textit{ビットコインの価格変動は、ファイナンシャルヘッジ（ステーブルコイン）、強い長期的な信念（ホドリング）、または何も理解せずに自分のお金の100%をビットコインに投資しないことなど、いくつかの解決策によって軽減することができます。 レビュー済み: true}

\item \textbf{Question 170:} C
\\ \textit{Bitcoinプロトコルがユニークなのは、世界規模で1日24時間、週7日間稼働しているため、金融当局による規制が難しいという点です。}

\item \textbf{Question 171:} B
\\ \textit{2017年は、特に何千もの初期コインオファリング（ICO）の開始とともに、暗号通貨界における顕著な投機バブルが起こった年でした。}

\item \textbf{Question 172:} C
\\ \textit{ビットコインは、2100万ビットコインという限定された数量と、平均して4年ごとに貨幣創造を半減させるプロセス（ハービング）により、一般的に上向きの傾向にあります。}

\item \textbf{Question 173:} B
\\ \textit{ビットコインが多くの投機的なバブルを経験するのは、実際に世界中どこでも使用される通貨であるためです。}

\item \textbf{Question 174:} C
\\ \textit{ビットコインプロトコルは、世界規模で1日24時間、週7日間稼働するという特徴を持っており、金融当局が規制を行うことを難しくしています。}

\item \textbf{Question 175:} D
\\ \textit{ビットコインは従来の市場にさらに統合されることで、生き残り、さらに成長を続けています。ビットコインETFの登場、より明確な規制、およびより良い保管ツールの取得は、この傾向を促進するだけです。}

\item \textbf{Question 176:} C
\\ \textit{Silk Roadのようなダークウェブマーケットプレイスでビットコインの使用が拡大したのは、その制御不可能で匿名性がある性質によるものです。}

\item \textbf{Question 177:} A
\\ \textit{ビットコインの価格は、顕著な変動性を特徴としており、市場の変動や強気市場と弱気市場の相場変動によって影響を受けることがあります。}

\item \textbf{Question 178:} C
\\ \textit{ビットコインは消えるバブルではなく、世界中どこでも実際に使用されている通貨です。 レビュー済み: true}

\item \textbf{Question 179:} C
\\ \textit{ビットコインは、フィアット通貨に対して並行経済と見なされており、これは伝統的な経済と並行して機能し、取引所プラットフォームを介さずに直接取引に使用できることを意味します。}

\item \textbf{Question 180:} A
\\ \textit{ビットコインを購入せずに入手する方法の一つは、提供する製品やサービスの支払い手段として受け入れることです。これにより、あなたの仕事を通じてビットコインを獲得することができます。}

\item \textbf{Question 181:} B
\\ \textit{マーチャントとしてビットコインを受け入れることには、検閲への抵抗、取引手数料の削減、効率の向上、インフレーションからの保護、 さらには財政的自由と主権など、いくつかの利点があります。 レビュー済み: true}

\item \textbf{Question 182:} B
\\ \textit{BTCMapは、Bitcoinを受け入れるすべての加盟店と世界中の異なるBitcoinコミュニティをリストアップするオープンソースで協力的なプロジェクトです。}

\item \textbf{Question 183:} A
\\ \textit{ビザンチン将軍問題は、分散コンピューティングにおける古典的な問題であり、一部のノードが故障しているか悪意がある可能性があるネットワークで合意に達することを含みます。ビットコインは、プルーフ・オブ・ワークに基づくブロックチェーンを使用することでこの問題を解決し、すべてのノードがトランザクション履歴の状態について合意することを保証します。}

\item \textbf{Question 184:} A
\\ \textit{AMF、または金融市場庁は、ビットコインが購入できるプラットフォームを規制するフランスの組織です。}

\item \textbf{Question 185:} C
\\ \textit{ビジネスでビットコインを受け入れるソリューションを選択する際には、予想される取引量、割り当てられた予算、ビジネスのタイプ（オンラインまたは実店舗）など、いくつかの要因を考慮する必要があります。 レビュー済み: true}

\item \textbf{Question 186:} B
\\ \textit{OpenNodeは、ビジネスがビットコインを受け入れるためのシンプルなオンラインソリューションです。他のソリューションには、アマチュア商人向けのSwiss Bitcoin Payや、大規模な構造または情熱的なビットコイナー向けのBTCpay Serverが含まれます。 レビュー済み: true}

\item \textbf{Question 187:} C
\\ \textit{BTCMapは、ビットコイン経済をよりアクセスしやすく、便利にするために貢献している取り組みの一つです。これは、ビットコインを受け入れるすべての商人と、世界中の異なるビットコインコミュニティをリストするオープンソースで協力的なプロジェクトです。 レビュー済み: true}

\item \textbf{Question 188:} C
\\ \textit{BTCpay Serverは、大規模な構造や情熱的なビットコイナーがBitcoinを受け入れるためのツールです。他の解決策には、シンプルなオンラインビジネス向けのOpenNodeや、アマチュア商人向けのSwiss Bitcoin Payが含まれます。
レビュー済み: true}

\item \textbf{Question 189:} D
\\ \textit{日常取引でビットコインの使用を容易にする一つの方法は、ビットコインを受け入れるすべての商人のリストを作成することです。これはBTCMapのようなプラットフォームを通じて行うことができます。 レビュー済み: true}

\item \textbf{Question 190:} C
\\ \textit{Swiss Bitcoin Payは、アマチュアの商人がビットコインを受け入れるためのソリューションです。他のソリューションには、シンプルなオンラインビジネス向けのOpenNodeや、大規模な構造や情熱的なビットコイナー向けのBTCpay Serverがあります。}

\item \textbf{Question 191:} B
\\ \textit{コースのテキストでは、ビットコインは非常にボラティリティが高い金融資産であり、その価格が0に下がる可能性があると述べられています。これはビットコインを購入する際に関連するリスクの一つです。}

\item \textbf{Question 192:} D
\\ \textit{コースのテキストでは、Bitcoinを購入する前に、鍵を保存する安全なウォレットを持っておくべきだと述べています。}

\item \textbf{Question 193:} B
\\ \textit{コースのテキストでは、ドルコスト平均法は定期的な間隔で少額のビットコインを購入する技術であると説明されています。この方法は、時間をかけて価格の変動を平準化し、所有するビットコインの量の継続的な成長を保証します。}

\item \textbf{Question 194:} B
\\ \textit{コースのテキストでは、自発的な購入はビットコインへの即時の露出を得るために使用されると説明されています。例えば、クラッシュ時に購入することや、ボーナスを利用することを意味することがあります。}

\item \textbf{Question 195:} A
\\ \textit{コースのテキストでは、「Know Your Customer」（KYC）規制が、テロ資金の調達、脱税、およびマネーロンダリングと戦うために、ユーザーに身分証明の提供を要求することが説明されています。}

\item \textbf{Question 196:} C
\\ \textit{コースのテキストには、KYCなしでビットコインを取得するいくつかの方法があると記載されており、その一つがビットコインATMです。}

\item \textbf{Question 197:} C
\\ \textit{コースのテキストでは、ビットコインを購入した後、ハッキングや資金ブロックのリスクを最小限に抑えるために、取引プラットフォームからUTXOを直ちに引き出す必要があるとアドバイスしています。}

\item \textbf{Question 198:} C
\\ \textit{このコースのテキストでは、DCAプラットフォームを選択する際、正しい選択はあなたの国での利用可能性により大きく依存すると述べています。
レビュー済み: true}

\item \textbf{Question 199:} C
\\ \textit{このテキストは、時々あなたのウォレット内のUTXOを統合することが、ビットコインを効果的に管理し、取引中の不必要な手数料を避けるために重要であると述べています。 レビュー済み: true}

\item \textbf{Question 200:} D
\\ \textit{このテキストは、FOMO（Fear of Missing Out：見逃し恐怖症）とFUD（Fear, Uncertainty, Doubt：恐怖、不確実性、疑念）が、ビットコインを購入する際に衝動的で潜在的に有害な意思決定につながる可能性のある感情であることを説明しています。}

\item \textbf{Question 201:} B
\\ \textit{このテキストでは、取引所プラットフォームにビットコインを残しておくことで、ハッキングと資金ブロックのリスクにさらされる可能性があると述べています。}

\item \textbf{Question 202:} D
\\ \textit{コースのテキストでは、あなたの国の規制に従わないでビットコインを購入すると、自分自身を危険にさらす可能性があると述べられています。}

\item \textbf{Question 203:} A
\\ \textit{ビットコインの採用は、他の新技術と同様に、Sカーブに従います。これは新技術の採用が一般的に見られるパターンで、採用はゆっくりと始まり、その後加速し、飽和に達すると減速します。}

\item \textbf{Question 204:} A
\\ \textit{エルサルバドルは大胆な一歩を踏み出し、ビットコインを完全に採用し、それを法定通貨として宣言しました。 レビュー済み: true}

\item \textbf{Question 205:} A
\\ \textit{ビットコインの台頭は、企業、大学、規制機関、個人にこの新しい技術を考慮に入れることを強いています。新しいツールを作成し、サービスを適応させ、生き残るためにはイノベーションを続けなければなりません。}

\item \textbf{Question 206:} A
\\ \textit{ビットコインは多分野にわたる科目として説明されています。これは、
暗号学、ゲーム理論、経済学と通貨政策、コンピュータサイエンス、哲学、エネルギー、法律、および規制に関連する問題を提起するためです。}

\item \textbf{Question 207:} D
\\ \textit{このコースのテキストは、ビットコイン革命の主な目標が、人間の主権が権利として認められ、プライバシーがデフォルトで尊重され、お金が操作されない世界を創造することであることを示唆しています。}

\item \textbf{Question 208:} B
\\ \textit{著名な経済学者であるミルトン・フリードマンは1999年に、信頼性の高い電子マネーが間もなく開発されるだろうと予測しました。これはインターネット上で、AがBを知らず、BがAを知らない状態で、AからBへ資金を移動できる方法です。}

\item \textbf{Question 209:} D
\\ \textit{Bitcoinの一般的な上昇傾向とは、ある期間にわたってその価格が一貫して増加していることを指します。上昇傾向はポジティブな軌道を示唆しているものの、価格が変動なしに上昇し続けることが保証されているわけではないことに注意することが重要です。}

\item \textbf{Question 210:} D
\\ \textit{コースのテキストでは、いくつかの国がビットコインの使用を禁止し、犯罪化したことが述べられており、文化、時代、国に基づいたビットコインの採用の複雑さが増しています。しかし、エルサルバドルはそれを法定通貨として採用する勇気を持ちました。}

\item \textbf{Question 211:} D
\\ \textit{ビットコインは、暗号学、ゲーム理論、経済学と金融政策、コンピュータ科学、哲学、エネルギー、法律、および規制を含む様々な分野に関連する疑問を提起する多分野にわたる主題です。}

\item \textbf{Question 212:} D
\\ \textit{コースのテキストでは、ビットコインの普及がバイラルであると説明されています。
これは、ビットコインが一人の人から別の人へと、ウイルスのように、しかし良い意味で、急速かつ広範囲に広がることを示唆しています。
レビュー済み: true}

\item \textbf{Question 213:} C
\\ \textit{コースのテキストは、ビットコインが消え去ることはないと示唆しています。それどころか、革命はちょうど始まったばかりです。}

\item \textbf{Question 214:} A
\\ \textit{このコースは、ビットコインには探求すべきことがたくさんあるため、一度にすべてを理解するのは複雑であることを示唆しています。そのため、時間をかけるべきだとしています。}

\item \textbf{Question 215:} A
\\ \textit{Lightning Networkは、Bitcoinプロトコルを使用して高速な取引を可能にする支払いネットワークです。別の暗号通貨ではなく、Bitcoinを採掘する方法でもありません。}

\item \textbf{Question 216:} C
\\ \textit{スケーラビリティ問題とは、採用されるにつれて、秒間の取引数を増やし続けることができる金融システムを実装するという課題を指します。}

\item \textbf{Question 217:} B
\\ \textit{ブロックチェーンのトリレンマとは、ブロックチェーンに基づくプロトコルが分散化、セキュリティ、スケーラビリティの基準のうち2つしか満たすことができないという課題を指します。}

\item \textbf{Question 218:} A
\\ \textit{Lightning Networkは即時で低コストのBitcoin取引を可能にし、Bitcoinのスケーラビリティ問題を解決します。 レビュー済み: true}

\item \textbf{Question 219:} A
\\ \textit{Lightning Networkは2017年に最初に検証され、実装されました。}

\item \textbf{Question 220:} A
\\ \textit{Lightning Networkは、送金者に対して即時取引を低料金で可能にする支払いチャネルのネットワークとして機能します。}

\item \textbf{Question 221:} A
\\ \textit{Western Union、中央銀行、Visa、Mastercardなどの従来の送金サービスは、Lightning Network技術を採用しない場合、消滅する可能性があります。}

\item \textbf{Question 222:} C
\\ \textit{Lightning Network上の取引手数料は非常に低く、多くの場合0.5%未満です。}

\item \textbf{Question 223:} B
\\ \textit{Lightning Network上のトランザクションは、暗号技術を通じて、および間接的にはBitcoin上のマイナーによって消費されるエネルギーを通じて保護されています。}

\item \textbf{Question 224:} A
\\ \textit{Lightning Networkは、直接のチャネル接続を持たない個人間での取引を可能にします。}

\item \textbf{Question 225:} D
\\ \textit{Bitcoinプロトコルにおいて、2つのブロックが生成される平均時間間隔は10分です。}

\item \textbf{Question 226:} C
\\ \textit{Bitcoinプロトコルにおけるブロックサイズは4MBです。}

\item \textbf{Question 227:} B
\\ \textit{Lightning Networkは、Bitcoinプロトコルの第二層ソリューションであり、非常に低い価値の取引を容易にすることを目的としています。これらの取引は、Bitcoinブロックチェーン上での高い手数料や長い確認時間のために、実用的でない場合があります。}

\item \textbf{Question 228:} A
\\ \textit{Lightning Networkは、ビットコインのセキュリティと分散化を確保するために必要な制約のために以前は手が届かなかった潜在的な応用の幅広い範囲を開きます。日常生活におけるこれらの応用の中で、コマースにおける即時請求を挙げることができます。 レビュー済み: true}

\item \textbf{Question 229:} B
\\ \textit{Lightning Networkは、ビットコインの小数単位であるサトシ（satsとも呼ばれます）を使用して機能します。}

\item \textbf{Question 230:} A
\\ \textit{「skin in the game」という概念は、最近ビデオゲーム業界で人気を博しているアイデアです。これは、開発者やステークホルダーがゲームの成功に個人的な投資や利害関係を持っているという考えを本質的に指し、それによって彼らの利益を品質の高い製品を提供することに合わせることを意味します。}

\item \textbf{Question 231:} B
\\ \textit{Lightning Networkは、プレイヤーがゲームをプレイする際に、数サトシといった非常に少額のお金を賭けることを可能にします。これにより、競争を促進する賭け金を設定することができ、同時にボットの展開コストを大幅に増加させます。}

\item \textbf{Question 232:} B
\\ \textit{Lightning Networkを使用すると、毎分マイクロトランザクションを行うことができ、消費者が消費する内容に基づいてコンテンツの支払いを行う経済モデルを実験する扉が開かれます。}

\item \textbf{Question 233:} B
\\ \textit{Lightning Networkのマイクロトランザクションでは、ユーザーは低い取引手数料を享受でき、非常に小さな金額のBitcoinを送金することが経済的になります。また、即時支払いが可能となり、リアルタイムでの取引が実現します。最後に、大量の取引を同時に処理できるスケーラビリティが得られます。したがって、誤った回答は長い支払い時間です。
確認済: true}

\item \textbf{Question 234:} D
\\ \textit{Lightning Networkを利用することで、コンテンツクリエーターはユーザーの注目を保持するために良質なコンテンツを制作するインセンティブを得ることができます。その結果、ユーザーは消費したコンテンツにのみ料金を支払うことになり、前払いのサブスクリプション料金が不要になります。}

\item \textbf{Question 235:} D
\\ \textit{DCAプラットフォームとは、「ドルコスト平均法」（DCA）を容易にするプラットフォームを指します。これは、個人が定期的な間隔で、資産の価格に関係なく、特定の資産に一定額のお金を投資する投資戦略です。}

\item \textbf{Question 236:} A
\\ \textit{Lightning Networkを利用することで、物品のレンタルにこのシステムを使用することが考えられます。ユーザーは、物品を使用した時間に応じて、分単位や秒単位で料金を請求されることがあります。}

\item \textbf{Question 237:} C
\\ \textit{Lightning Networkは、Bitcoinユーザーにとって多くの興味深い使用例を開放します。その結果として生まれる経済モデルやビジネスチャンスは数多く、消費した内容に基づいて消費者が支払う新しい経済モデルを含んでいます。}

\item \textbf{Question 238:} A
\\ \textit{Lightning Networkを自分で試す一つの方法は、お気に入りのポッドキャストを聴くことで少量のsats（サトシ）を報酬として得られるポッドキャストアプリケーション「Fountain」を試すことです。 レビュー済み: true}

\item \textbf{Question 239:} A
\\ \textit{テキストには、AI技術の進歩の一つの潜在的な影響として、AIと自動化により80％以上の仕事が消滅する可能性があると記載されています。}

\item \textbf{Question 240:} C
\\ \textit{このテキストは、ビットコインが通信手段のインターネットと同様に、技術革命であり、信頼できる第三者なしに価値を交換する新しい大規模組織の形態を生み出す可能性を私たちにもたらすことを説明しています。}

\item \textbf{Question 241:} B
\\ \textit{このテキストは金融の未来についていくつかの疑問を提起しており、その一つが私たちが使用するお金を誰が保持し、認証し、追跡すべきかということです。}

\item \textbf{Question 242:} B
\\ \textit{テキストには、Bitcoinが擬似匿名であり、どのような権限からも認可を受けることなく世界中どこでも交換できると述べられています。}

\item \textbf{Question 243:} A
\\ \textit{このテキストは、これらの新しい技術を受け入れることが世界中で大規模な規模の経済を生み出す可能性があることを示唆しています。 レビュー済み: true}

\item \textbf{Question 244:} A
\\ \textit{テキストでは、銀行システムを誰がコントロールすべきかという問題が重要であると述べられています。なぜなら、銀行のルールが透明ではなく、すべての人にとって理解しやすいものではないからです。}

\item \textbf{Question 245:} A
\\ \textit{このテキストは、検閲を容認することが表現の自由と集会の権利を破壊する可能性があると述べています。 レビュー済み: true}

\item \textbf{Question 246:} C
\\ \textit{テキストには、世界には銀行口座を持たない24億人の人々がいると述べられており、ビットコインは彼らの社会的地位や政治的立場に関係なく、取引において平等を実現します。}

\item \textbf{Question 247:} C
\\ \textit{このテキストでは、Bitcoinプロトコルがすべてのユーザーに対して中立であると述べられています。}

\item \textbf{Question 248:} A
\\ \textit{このテキストでは、ユーザーがビットコインを取引で受け入れるか、規制されたプラットフォームまたは規制されていないプラットフォームを通じて購入することでビットコインを入手できると述べています。
レビュー済み: true}

\item \textbf{Question 249:} C
\\ \textit{テキストには、サトシ・ナカモトが2008年に金融システムの変更を提案するためにビットコインを作成したと記載されています。}

\item \textbf{Question 250:} C
\\ \textit{このテキストでは、ビットコインが銀行システムからの自立を可能にすると述べられています。}


\end{enumerate}

\end{document}
